\documentclass[11pt]{article}
\usepackage[T1]{fontenc}
\usepackage[utf8]{inputenc}
\usepackage{geometry}
\usepackage{lmodern}
\usepackage{microtype}
\usepackage{amsmath,amssymb,amsthm}
\usepackage{hyperref}
\usepackage{booktabs}
\usepackage{enumitem}
\usepackage{xurl}
\geometry{margin=1in}

\newtheorem{theorem}{Theorem}
\newtheorem{lemma}[theorem]{Lemma}
\newtheorem{proposition}[theorem]{Proposition}
\newtheorem{corollary}[theorem]{Corollary}
\theoremstyle{definition}
\newtheorem{definition}[theorem]{Definition}
\newtheorem{hypothetical}[theorem]{Hypothetical Criterion}
\theoremstyle{remark}
\newtheorem{remark}[theorem]{Remark}

\title{Certified Log-Concavity of the Riemann--Jacobi Kernel\\
and a Source Audit of a P\'{o}lya-Type Real-Zero Criterion}
\author{Tristen Kyle Pierson\\
  Axiovex Systems\\
  \small ORCID: \href{https://orcid.org/0009-0003-7269-956X}{0009-0003-7269-956X}}
\date{June 2026; revised November 2026}

% Version DOI of the previous version (v2). Update to the v3 DOI at deposit.
\newcommand{\doiurl}{\url{https://doi.org/10.5281/zenodo.24445687}}
\newcommand{\repourl}{\url{https://github.com/AXIOVEX/riemann-solver}}

\begin{document}
\maketitle

\begin{abstract}
We certify strict log-concavity of the Riemann--Jacobi kernel~$\Phi(u)$
appearing in the Fourier cosine representation
\[
\Xi(t) = \int_0^\infty \Phi(u)\cos(tu)\,du
\]
of the Riemann Xi function.
The proof combines an elementary algebraic analysis of the dominant term,
rigorous interval arithmetic on~$[0,1]$, a no-cancellation certification
of $(\log\Phi)'' < 0$ on~$[1,3]$, and an explicit monotonicity tail bound for~$u\geq 3$.

The strict log-concavity of the full kernel~$\Phi$ was first proved by
Csordas~\cite{csordas2015} (Theorem~4.5); the certification given here is
independent of that proof.

We also audit the commonly cited P\'{o}lya-type bridge from log-concavity of a
kernel to real-rootedness of its Fourier transform.
Direct inspection of P\'{o}lya (1927) Satz~I and Satz~II shows that these theorems
are, respectively, a preservation theorem and a Mellin--Fourier theorem; neither states
the log-concavity criterion needed to derive the Riemann Hypothesis from the result
proved here.
More decisively, the criterion is false in general: the kernel
$K(t) = (5+t^4)e^{-t^2-t^4/1000}$ satisfies all five of its hypotheses, but its
Fourier transform has non-real zeros (Proposition~\ref{prop:counterexample}).
Accordingly, \textbf{no unconditional RH claim is made}, and the conditional
route through the criterion is closed: the criterion cannot be established,
because it is not true.
This source audit is part of the contribution: it separates the certified
log-concavity theorem from an often-assumed real-rootedness criterion, and
settles the status of that criterion.

The unconditional contributions of this paper are the independent
interval-arithmetic certification of the strict log-concavity
of~$\Phi$ on~$[0,\infty)$ and the source audit, including the refutation of
the criterion.
The proof of log-concavity is independent of the P\'{o}lya-source audit.

\medskip\noindent
\textbf{DOI:} \doiurl
\end{abstract}

%% =========================================================================
\section{Introduction}
\label{sec:intro}

The Riemann Hypothesis (RH), posed by Bernhard Riemann in 1859, asserts that
all nontrivial zeros of the Riemann zeta function~$\zeta(s)$ lie on the
critical line~$\operatorname{Re}(s) = 1/2$.
Equivalently, all zeros of the Riemann Xi function
\begin{equation}\label{eq:xi}
\Xi(t) := \xi\!\left(\tfrac{1}{2} + it\right), \quad
\xi(s) := \tfrac{1}{2}s(s-1)\pi^{-s/2}\Gamma(s/2)\zeta(s),
\end{equation}
are real.
It is well known~\cite{titchmarsh1986} that~$\Xi(t)$ admits the Fourier cosine
representation
\begin{equation}\label{eq:fourier}
\Xi(t) = \int_0^\infty \Phi(u)\cos(tu)\,du,
\end{equation}
where the \emph{Riemann--Jacobi kernel} is
\begin{equation}\label{eq:phi}
\Phi(u) = 4\sum_{n=1}^{\infty}\varphi_n(u), \quad
\varphi_n(u) = \bigl(2\pi^2 n^4 e^{9u/2} - 3\pi n^2 e^{5u/2}\bigr)e^{-\pi n^2 e^{2u}}.
\end{equation}
We verify computationally that
$\int_0^\infty \Phi(u)\,du = \xi(1/2) \approx 0.4971$ to 15~digits.

\paragraph{Approach.}
This paper proves a structural property of the Riemann--Jacobi kernel:
strict log-concavity of~$\Phi$ on~$[0,\infty)$ (Theorem~\ref{thm:main}).
Because the Fourier transform of this kernel is~$2\Xi$, this property is
naturally relevant to the real-zero problem for~$\Xi$.
However, the implication from kernel log-concavity to real-rootedness of the
Fourier transform is not established by the P\'{o}lya sources audited here.
We therefore separate the unconditional log-concavity theorem from the conditional
RH implication.

The paper should be read as a certification of log-concavity, not as a proof of RH.
The RH consequence appears only in a conditional remark
(Remark~\ref{thm:rh_conditional}), dependent on a hypothetical criterion that
is false (Proposition~\ref{prop:counterexample}); the conditional statement is
therefore vacuous.
The source audit (Section~\ref{sec:polya_audit_main}) is not needed for the proof
of Theorem~\ref{thm:main}; it only governs the interpretation of possible RH
consequences.

A second contribution of the paper is therefore source-critical: we trace the commonly
cited log-concavity-to-real-rootedness bridge through P\'{o}lya (1927),
Csordas--Varga (1989), de~Bruijn (1950), Ilieff (1955), and later literature, and
find no verified primary-source statement matching
Hypothetical Criterion~\ref{hcrit:polya_star}.
The audit does not stop at the missing source: in
Section~\ref{sec:polya} we prove the criterion is false, by an explicit
counterexample satisfying all five of its hypotheses.
In this paper the criterion is consequently treated as a refuted
dependency, not as an established theorem and not as an open problem.

\paragraph{Prior work.}
The log-concavity of the dominant term~$\varphi_1$ was observed by Csordas and
Varga~\cite{csordas1989}.
The Tur\'{a}n inequalities for the Taylor coefficients of~$\Xi$ were established
by Csordas, Norfolk, and Varga~\cite{csordas1986}, and higher-order generalizations
by Dimitrov and Lucas~\cite{dimitrov2011} and Griffin, Ono, Rolen, and
Zagier~\cite{griffin2019}.
The connection between log-concavity and the de~Bruijn--Newman constant was explored
by de~Bruijn~\cite{debruijn1950} and Rodgers and Tao~\cite{rodgers2020}.
Csordas~\cite{csordas2015} proved the strict log-concavity of the
\emph{full kernel}~$\Phi$ (Theorem~4.5), via the stronger statement that
$\log\Phi(\sqrt{t})$ is concave in~$t$ (Theorem~4.2(b) of~\cite{csordas2015},
after Csordas--Varga (1988)).
The contribution of the present work is therefore not the log-concavity
theorem itself, but an independent interval-arithmetic certification of it,
with every estimate derived in Appendix~\ref{sec:appE} and every certificate
reproducible from the repository, together with the source audit and the
refutation of the criterion.

\paragraph{Claims intentionally not made.}
This paper does not claim:
\begin{enumerate}[leftmargin=2em,itemsep=1pt]
\item that RH is proved;
\item that P\'{o}lya (1927) Satz~I or Satz~II implies the log-concavity criterion;
\item that Csordas--Varga (1989) Theorem~2.2 states the log-concavity criterion;
\item that Lean verifies RH;
\item that the interval arithmetic certificates are formalized in Lean;
\item that Hypothetical Criterion~\ref{hcrit:polya_star} is established in the
peer-reviewed literature (it is false; Proposition~\ref{prop:counterexample});
\item that the strict log-concavity of~$\Phi$ is new (it is
Csordas's theorem~\cite[Theorem~4.5]{csordas2015}; the certification here
is independent).
\end{enumerate}

%% =========================================================================
\section{Properties of $\Phi$}
\label{sec:phi}

\begin{proposition}\label{prop:phi_properties}
The kernel~$\Phi$ defined by~\eqref{eq:phi} satisfies:
\begin{enumerate}[label=(\roman*)]
\item $\Phi(u) > 0$ for all~$u \geq 0$;
\item $\Phi(-u) = \Phi(u)$ for all~$u$;
\item $\Phi \in L^1(\mathbb{R})$;
\item $\Phi(u) = O\bigl(e^{-(\pi-\eta) e^{2u}}\bigr)$ as~$u \to +\infty$,
  for each $0 < \eta < \pi$;
\item $\Phi$ is real analytic on~$\mathbb{R}$.
\end{enumerate}
\end{proposition}

\begin{proof}
Properties~(i)--(iii) are classical; see~\cite[\S2.10]{titchmarsh1986}
and~\cite[Theorem~A]{csordas1989}.
For~(iv): the dominant term satisfies
$\varphi_1(u) \sim 2\pi^2 e^{9u/2} e^{-\pi e^{2u}}$ as $u \to +\infty$
(the prefactor $2\pi^2 e^{9u/2}$ grows, so the stronger statement
$\Phi(u) = O(e^{-\pi e^{2u}})$ printed in earlier versions is false);
hence $\Phi(u) \sim 8\pi^2 e^{9u/2} e^{-\pi e^{2u}}$, and for each
$0 < \eta < \pi$ the growing prefactor is absorbed by
$e^{-\eta e^{2u}}$, giving $\Phi(u) = O(e^{-(\pi-\eta)e^{2u}})$.
For~(v): fix any $u_0\in\mathbb{R}$.
For $r = \pi/8$, any $z$ with $|z - u_0| < r$ satisfies
$\operatorname{Re}(e^{2z}) \geq e^{2(u_0-r)}/\sqrt{2} =: c > 0$.
Thus $|\varphi_n(z)| \leq A\,n^4\,e^{-\pi c n^2/2}$ uniformly on $|z-u_0|\leq r/2$.
Since $\sum n^4 e^{-\pi c n^2/2} < \infty$, the Weierstrass $M$-test gives uniform
convergence; hence $\Phi$ is real analytic at $u_0$.
($\Phi$ is \emph{not} entire; only real analyticity on~$\mathbb{R}$, as required by
condition~(v) of Hypothetical Criterion~\ref{hcrit:polya_star}, is claimed.)

\emph{Numerical verification.}
$\Phi(u) > 0$ at 10{,}001 uniformly spaced points on $[0,1]$
($\min = 5.51 \times 10^{-7}$ at $u = 1$);
$|\Phi(u) - \Phi(-u)|/|\Phi(u)| < 10^{-70}$ at 8~test points.
\end{proof}

%% =========================================================================
\section{Algebraic Core}
\label{sec:algebraic}

\begin{theorem}[Algebraic Core]\label{thm:algebraic}
$(\log\varphi_1)''(u) < 0$ for all~$u \geq 0$.
\end{theorem}

\begin{proof}
Write $\varphi_1 = g \cdot E$ where $g(u) = \pi e^{5u/2}h(u)$,
$h(u) = 2\pi e^{2u} - 3$, and $E(u) = e^{-\pi e^{2u}}$.
Since $h(u) \geq 2\pi - 3 > 0$ for~$u \geq 0$, we have $\varphi_1 > 0$. Then
\[
\log\varphi_1 = \log\pi + \tfrac{5}{2}u + \log h(u) - \pi e^{2u},
\quad
(\log\varphi_1)'' = (\log h)'' - 4\pi e^{2u}.
\]
With $h' = 4\pi e^{2u}$ and $h'' = 8\pi e^{2u}$:
\[
(\log h)'' = \frac{h''h - (h')^2}{h^2} = \frac{-24\pi e^{2u}}{h(u)^2} < 0.
\]
Both terms are negative, so $(\log\varphi_1)'' < 0$.
\end{proof}

\begin{remark}
$(\log\varphi_1)''(0) \approx -19.56$, monotonically decreasing.
The series for~$\Phi''(0)$ converges to $-68.052\ldots$ by $N = 5$ terms.
\end{remark}

%% =========================================================================
\section{Tail Estimate}
\label{sec:tail}

\begin{lemma}\label{lem:tail_decay}
For $n \geq 2$ and $u \geq 0$: $e^{-\pi n^2 e^{2u}} \leq e^{-3\pi}\cdot e^{-\pi e^{2u}}$.
\end{lemma}

\begin{proof}
Equivalent to $(n^2 - 1)e^{2u} \geq 3$, holding since $n^2 - 1 \geq 3$ and $e^{2u} \geq 1$.
\end{proof}

\begin{proposition}\label{prop:tail_ratio}
For $u \geq 0$: $|R(u)|/\varphi_1(u) < 1/50$, where $R = \sum_{n\geq 2}\varphi_n$.
\end{proposition}

\begin{proof}
$|\varphi_n|/\varphi_1 \leq 2n^4 e^{-\pi(n^2-1)e^{2u}}$ (since $B_n(u)/n^4 < 2$).
Each majorant term $2n^4 e^{-\pi(n^2-1)e^{2u}}$ is decreasing in~$u$
(the exponent $-\pi(n^2-1)e^{2u}$ decreases as $e^{2u}$ grows), so the
sum over $n \geq 2$ is decreasing in~$u$ and its value at $u = 0$ bounds
all $u \geq 0$:
at $u = 0$: $\sum_{n\geq 2} 2n^4 e^{-\pi(n^2-1)} < 0.006 < 1/50$.
\end{proof}

\begin{remark}\label{rem:truncation}
The interval arithmetic retains $n = 1,\ldots,5$ and omits
$\delta := 4\sum_{n \geq 6}\varphi_n$.
Uniformly on~$[0,1]$, interval arithmetic certifies
(\texttt{proof/verify\_truncation\_and\_crosscheck.py}):
$\sup|\delta| \leq 7.13 \times 10^{-43}$,
$\sup|\delta'| \leq 1.20 \times 10^{-39}$,
$\sup|\delta''| \leq 2.01 \times 10^{-36}$,
from explicit majorants on $|\varphi_n^{(k)}|$ for $n \geq 6$
(growth factors at $u=1$, decay $e^{-\pi n^2}$ at $u=0$).
Propagating through
$Q_\Phi - Q_{\Phi_5} = \Phi_5\delta'' + \Phi_5''\delta
- 2\Phi_5'\delta' + \delta\delta'' - (\delta')^2$
with the certified enclosures
$\sup|\Phi_5| \leq 2.65$, $\sup|\Phi_5'| \leq 13.9$,
$\sup|\Phi_5''| \leq 143$ on~$[0,1]$ gives the uniform bound
$|Q_\Phi - Q_{\Phi_5}| \leq 5.34 \times 10^{-36}$ on~$[0,1]$.
Since the certified margin is $|Q_{\Phi_5}| \geq 3.36 \times 10^{-12}$,
the safety factor exceeds $6 \times 10^{23}$.
(Earlier versions of this remark propagated the cross terms only at
$u=1$ and printed a safety factor of $2.9 \times 10^{30}$; that
derivation was not uniform on~$[0,1]$ and the figure is withdrawn.)
\end{remark}

%% =========================================================================
\section{Interval Arithmetic Verification}
\label{sec:ia}

\begin{definition}
The \emph{log-concavity numerator} of a positive function~$f$ is
$Q_f(u) := f''(u)f(u) - f'(u)^2$.
\end{definition}

\begin{theorem}\label{thm:ia}
$Q_\Phi(u) < 0$ for all~$u \in [0, 1]$.
\end{theorem}

\begin{proof}
Closed-form derivatives:
\begin{align*}
g_n' &= 9\pi^2 n^4 e^{9u/2} - \tfrac{15}{2}\pi n^2 e^{5u/2}, &
g_n'' &= \tfrac{81}{2}\pi^2 n^4 e^{9u/2} - \tfrac{75}{4}\pi n^2 e^{5u/2}, \\
E_n' &= -2\pi n^2 e^{2u}\cdot E_n, &
E_n'' &= \bigl(-4\pi n^2 e^{2u} + 4\pi^2 n^4 e^{4u}\bigr)\cdot E_n.
\end{align*}
The primary certificate is \texttt{proof/verify\_logconcavity\_arb.py}, using
Arb/FLINT ball arithmetic.
The independent cross-check is \texttt{proof/verify\_logconcavity\_rigorous.py}, using
\texttt{mpmath.iv} interval arithmetic at 60-digit precision.
Both retain $n = 1,\ldots,5$ and use the same closed-form derivative formulas above.

The \texttt{mpmath.iv} cross-check partitions $[0, 0.949]$ into 1{,}898 subintervals
and $[0.949, 1.0]$ into 51{,}000 subintervals.
The Arb/FLINT verifier uses its own interval partition and precision settings,
reported in \texttt{results/verify\_logconcavity\_arb.json}.
On every interval, the verifier evaluates~$Q_{\Phi_5}$ by interval or ball
arithmetic and certifies that the upper bound is negative.

\textbf{Result:} The Arb/FLINT primary verifier certifies $Q_{\Phi_5}(u) < 0$
on $[0,1]$, and the \texttt{mpmath.iv} cross-check reproduces the sign
certification on all 52{,}898 subintervals, with maximum \texttt{mpmath.iv}
upper bound $-3.36 \times 10^{-12}$.
Since the uniform truncation error is $\leq 5.34 \times 10^{-36}$
(Remark~\ref{rem:truncation}),
$Q_\Phi(u) < 0$ on~$[0,1]$.
\end{proof}

%% =========================================================================
\section{Extended Certification on $[1.0, 3.0]$}
\label{sec:extended}

Direct interval arithmetic on $Q_\Phi$ fails for $u > 1$ due to catastrophic
cancellation ($40\times$).
We certify $(\log\Phi)''(u) < 0$ directly.

\begin{lemma}[Uniform tail bound on {$[1,3]$}]\label{lem:wtail}
For all $u\in[1,3]$, $|W_{\mathrm{tail}}(u)| \leq B(u)$,
where $W_{\mathrm{tail}}(u) := (\log\Phi)''(u) - (\log\varphi_1)''(u)$
and $B(u)$ is the interval-arithmetic upper bound defined in Appendix~\ref{sec:appE}.
The verifier certifies $B(u) \leq 8.54\times 10^{-24}$ on the full cover of $[1,3]$.
The sign proof uses the Arb ball comparison $W_1(I)+B(I)<0$ on each ball;
the scalar bound $B\leq 8.54\times 10^{-24}$ is a diagnostic summary.

Complete derivation in Appendix~\ref{sec:appE}.
\end{lemma}

\begin{theorem}\label{thm:extended}
$(\log\Phi)''(u) < 0$ for all~$u \in [1.0, 3.0]$.

All proof-critical estimates used in this certification are derived in
Appendix~\ref{sec:appE}; the scripts are reproducibility checks and certify
numerical constants, not substitutes for the written proof.
\end{theorem}

\begin{proof}
Decompose $(\log\Phi)'' = W_1 + W_{\mathrm{tail}}$ where
$W_1(u) = -24\pi e^{2u}/h(u)^2 - 4\pi e^{2u}$ is proved negative in
Theorem~\ref{thm:algebraic}.
The verifier
(\texttt{proof/verify\_ia\_1\_to\_3\_arb.py}, Appendix~\ref{sec:appE})
certifies by Arb/FLINT ball arithmetic that
$W_1(I) + B(I) < 0$ (a ball comparison in Arb) on each of 1{,}813
adaptive balls~$I$ whose union is $[0.99,3.01]\supset[1.0,3.0]$;
the half-width at left endpoint~$u$ is
$\min(0.01,\, 0.35/(2\pi e^{2u}))$, so the exponent $-\pi e^{2u}$ of the
dominant term varies by at most $0.35$ across any one ball.
Here $B(u)$ is the upper bound on $|W_{\mathrm{tail}}(u)|$
of Appendix~\ref{sec:appE}, Lemma~\ref{lem:Bsign}, computed from a
direct enclosure of~$\Delta Q$.

\textbf{Result:} All 1{,}813 checks certified; maximum
$B \leq 8.54\times 10^{-24}$;
minimum margin $(-W_1 - B) = 91.29$ at $u\approx1.0$.
Since $|W_{\mathrm{tail}}|\leq B$:
$(\log\Phi)'' = W_1 + W_{\mathrm{tail}} \leq W_1 + B < 0$.
\end{proof}

%% =========================================================================
\section{Tail Bound for $u \geq 3.0$}
\label{sec:tailbound}

\begin{theorem}\label{thm:perturbation_tail}
$(\log\Phi)''(u) < 0$ for all~$u \geq 3.0$.
\end{theorem}

\begin{proof}
For $u\ge3$, Appendix~\ref{sec:appE7}
(Lemmas~\ref{lem:weighted_comparisons}--\ref{lem:fartail_mono}) proves
\[
|W_{\mathrm{tail}}(u)|\le C_{\mathrm{tail}}\,e^{4u}\varepsilon^*(u),
\]
where $C_{\mathrm{tail}}=2600$.
This is entirely separate from and independent of the interval-certified $B(u)$
bound used on~$[1,3]$.

At $u=3$,
\[
W_1(3)<-4\pi e^6<-1000,\qquad
C_{\mathrm{tail}}\,e^{12}\,\varepsilon^*(3)
\leq 2600\,e^{12}\cdot 32\,e^{-3\pi e^6}
< 10^{-1629} < 1.
\]
Moreover, $W_1(u)$ is strictly decreasing on $[3,\infty)$
(Lemma~\ref{lem:W1_mono}) and $e^{4u}\varepsilon^*(u)$ is strictly decreasing on
$[3,\infty)$ (Lemma~\ref{lem:fartail_mono}).
Therefore
\[
W_1(u)+C_{\mathrm{tail}}\,e^{4u}\varepsilon^*(u)
\le
W_1(3)+C_{\mathrm{tail}}\,e^{12}\varepsilon^*(3)<0.
\]
Since $|W_{\mathrm{tail}}|\le C_{\mathrm{tail}}\,e^{4u}\varepsilon^*$,
\[
(\log\Phi)''(u)=W_1(u)+W_{\mathrm{tail}}(u)<0.
\]

Table: $\varepsilon^*(1.0) \approx 1.82\times 10^{-29}$,
$\varepsilon^*(1.5) \approx 1.96\times 10^{-81}$,
$\varepsilon^*(2.0) \approx 1.07\times 10^{-222}$,
$\varepsilon^*(3.0) \leq 1.66\times 10^{-1650}$ (certified,
\texttt{proof/verify\_weighted\_comparisons.py}).
\end{proof}

%% =========================================================================
\section{Main Result}
\label{sec:main}

\begin{theorem}[Log-concavity of the Riemann--Jacobi kernel]\label{thm:main}
The kernel~$\Phi$ defined by~\eqref{eq:phi} satisfies $(\log\Phi)''(u) < 0$
for all~$u \geq 0$.
\end{theorem}

\begin{proof}
Combine Theorems~\ref{thm:ia} ($u \in [0,1]$), \ref{thm:extended} ($u \in [1.0, 3.0]$),
and~\ref{thm:perturbation_tail} ($u \geq 3.0$).
\end{proof}

%% =========================================================================
\section{Source-Critical Audit of P\'{o}lya (1927)}
\label{sec:polya_audit_main}

The theorem bridge between log-concavity of~$\Phi$ and real-rootedness of~$\Xi$ is
attributed in the literature to P\'{o}lya (1927)~\cite{polya1927}.
This section is \textbf{not part of the proof} of Theorem~\ref{thm:main}.
It audits whether the commonly cited P\'{o}lya sources support the log-concavity
criterion, and establishes the source-dependency status for the conditional
RH consequence.

\medskip
\noindent\textbf{Source.}
G.\ P\'{o}lya, ``\"{U}ber trigonometrische Integrale mit nur reellen Nullstellen,''
\textit{J.\ reine angew.\ Math.}\ \textbf{158} (1927), pp.~6--18.
Digitisation: Niedersächsische Staats- und Universitätsbibliothek Göttingen (GDZ),
LOG\,Id: LOG\_0006.
Literal title translation: ``On trigonometric integrals with only real zeros.''

\subsection{Exact German text of Satz~I (p.~7)}\label{sec:satzi}

\begin{quote}\itshape
I.\ Die Funktion $\varphi(t)$ sei analytisch f\"{u}r alle reellen Werte von $t$
und von folgender Beschaffenheit: Wenn das Integral
\[\int_{-\infty}^{+\infty} F(t)\,e^{izt}\,dt\]
nur reelle Nullstellen hat, so hat (wie auch $F(t)$ sonst beschaffen sei) das Integral
\[\int_{-\infty}^{+\infty} \varphi(t)\,F(t)\,e^{izt}\,dt\]
ebenfalls nur reelle Nullstellen.  Hierzu ist notwendig und hinreichend,
da\ss{} $C\,e^{\gamma^{2}}\varphi(it)$ eine ganze Funktion vom Geschlecht~1 sei,
die nur reelle Nullstellen besitzt und f\"{u}r reelles $t$ reelle Werte annimmt;
$C$ bedeutet eine nichtverschwindende, $\gamma$ eine reelle, nichtnegative Konstante.
\end{quote}

\paragraph{Literal English translation of Satz~I.}
\begin{quote}
I.\ The function $\varphi(t)$ shall be analytic for all real values of~$t$ and of
the following nature: if the integral
$\int_{-\infty}^{+\infty} F(t)\,e^{izt}\,dt$ has only real zeros,
then (whatever the nature of~$F(t)$ otherwise) the integral
$\int_{-\infty}^{+\infty} \varphi(t)\,F(t)\,e^{izt}\,dt$ also has only real zeros.
For this, the necessary and sufficient condition is that
$C\,e^{\gamma^{2}}\varphi(it)$ be an entire function of genus~$1$
possessing only real zeros and taking real values for real~$t$;
$C$ denotes a non-vanishing constant, $\gamma$ a real, non-negative constant.
\end{quote}

\paragraph{Assessment.}
In modern terms, the universal multipliers are, up to the allowed
nonzero constant, those whose rotation $\varphi(i\,\cdot)$ belongs to
the Laguerre--P\'{o}lya class~$\mathcal{LP}$ (the rotation is essential:
$e^{t^2}$ is a universal multiplier but does not itself belong
to~$\mathcal{LP}$).
Satz~I is a \textbf{preservation theorem}: it requires the Fourier transform of the
input kernel to already have only real zeros---which is precisely the conclusion
sought for~$\Xi$---and so cannot establish real-rootedness from scratch.

\subsection{Exact German text of Satz~II (pp.~7--8)}\label{sec:satzii}

\begin{quote}\itshape
II.\ Es sei die Funktion $f(t)$ f\"{u}r positive Werte von $t$ definiert,
reellwertig, absolut integrabel und f\"{u}r gen\"{u}gend gro\ss{}es $t$
der Ungleichung
\[|f(t)| < B\,e^{-t^{1+\beta}}\]
unterworfen ($B$, $\beta$ sind positive Konstanten).
Ferner sei $f(t)$ in einer Umgebung des Punktes $t = 0$ analytisch.
Dann stellt das Integral
\[\int_0^{\infty} t^{z-1}\,f(t)\,dt = H(z)\]
eine meromorphe Funktion dar.
Wenn die Funktion $H(z)$ nur reelle negative Nullstellen besitzt,
und $q$ eine positive ganze Zahl ist, so besitzt die durch das Integral
\[\int_{-\infty}^{+\infty} f(t^{2q})\,e^{izt}\,dt\]
dargestellte ganze Funktion nur reelle Nullstellen.
\end{quote}

\paragraph{Literal English translation of Satz~II.}
\begin{quote}
II.\ Let the function $f(t)$ be defined for positive values of~$t$, real-valued,
absolutely integrable, and for sufficiently large~$t$ subject to the inequality
$|f(t)| < B\,e^{-t^{1+\beta}}$ ($B$,~$\beta$ positive constants).
Furthermore let $f(t)$ be analytic in a neighbourhood of the point $t = 0$.
Then the integral $\int_0^{\infty} t^{z-1} f(t)\,dt = H(z)$ represents a
meromorphic function.
If the function $H(z)$ has only real negative zeros, and $q$ is a positive
integer, then the entire function represented by the integral
$\int_{-\infty}^{+\infty} f(t^{2q})\,e^{izt}\,dt$ has only real zeros.
\end{quote}

\paragraph{Assessment.}
The key hypothesis is that $H(z)$ has only real negative zeros---a condition on the
\textbf{Mellin transform} $H(z)=\int_0^\infty t^{z-1}f(t)\,dt$, not a log-concavity
condition on~$f$.
Log-concavity of~$f$ does not appear anywhere in Satz~II.
The conclusion concerns the Fourier integral of $f(t^{2q})$, not of~$f$ directly.

\subsection{Hypothesis comparison}\label{sec:interp_table}

\begin{center}\small
\begin{tabular}{@{}p{3.0cm}p{4.8cm}p{2.5cm}p{2.0cm}@{}}
\toprule
Interpretation & Theorem form &
Implies RH from log-concavity of~$\Phi$? & Status \\
\midrule
Strong log-concavity criterion &
  Positivity + log-concavity + analyticity + decay $\Rightarrow$
  Fourier transform has only real zeros &
  Yes & \textbf{Not found in P\'{o}lya 1927} \\[6pt]
Satz~I: Preservation theorem &
  IF $\int Fe^{izt}dt$ has real zeros, THEN multiplying by $\varphi\in\mathcal{LP}$
  preserves real zeros &
  No (circular) & Confirmed from source \\[6pt]
Satz~II: Mellin--Fourier theorem &
  IF Mellin transform $H(z)$ has only real negative zeros THEN
  $\int f(t^{2q})e^{izt}dt$ has only real zeros &
  No (different hypothesis) & Confirmed from source \\
\bottomrule
\end{tabular}
\end{center}

\subsection{Audit scope and verdict}\label{sec:verdict}

\begin{center}
\fbox{\parbox{0.93\textwidth}{%
\textbf{Verdict D --- source mismatch.}\\[3pt]
Direct reading of Satz~I and Satz~II of P\'{o}lya (1927)~\cite{polya1927} confirms
neither contains the log-concavity criterion.
\begin{itemize}[leftmargin=1.4em,topsep=3pt,itemsep=1pt]
\item Satz~I: preservation theorem; directly read and confirmed as mismatch.
\item Satz~II: Mellin--Fourier theorem; log-concavity does not appear;
  directly read and confirmed as mismatch.
\item Hilfss\"{a}tze (auxiliary lemmas): The present audit
  transcribes and analyzes Satz~I and Satz~II of P\'{o}lya (1927) only.
  The one auxiliary lemma that has been cited elsewhere under the name
  ``P\'{o}lya's Hilfssatz~II'' is now identified: it is Lemma~2 of
  Cardon~\cite{cardon2005}, who sources it to P\'{o}lya (1926)~\cite{polya1926},
  not to the 1927 Crelle paper.
  Hilfssatz~II is a shift-sum preservation lemma (for $a > 0$ and $b$ real,
  if $G$ is a real entire function of genus~0 or~1 with only real zeros,
  then $e^{ib}G(z+ia) + e^{-ib}G(z-ia)$ has only real zeros); it is not
  the log-concavity criterion, and the verification caveat carried in
  earlier versions of this audit is closed.
  (An earlier draft attributed the Hilfssatz~II discussion to
  Cardon (2002)~\cite{cardon2002}; that is a distinct paper, on
  convolution operators.)
\end{itemize}
\textbf{Consequence:}
Peer-reviewed literature (Coffey--Csordas 2013~\cite{csordas2013};
Csordas 2015~\cite{csordas2015},
Theorems~3.5--3.7 and introduction, who states ``no known necessary and sufficient
conditions'') indicates this criterion is not a known theorem.
Csordas (2015), Section~3, proves that within the admissible-kernel framework,
strict log-concavity is not by itself the known sufficient condition for
real-rootedness; additional generalized Laguerre inequalities enter the criterion
(Csordas 2015, Theorem~3.5).
A 2026 preprint~\cite{gershon2026} attributes a simple version to P\'{o}lya (1927)
but this cannot be verified against the primary text and is not supported by the
peer-reviewed record.
This paper settles the question the audit opened: the criterion is
\textbf{false} (Proposition~\ref{prop:counterexample}).
The RH conditional consequence (Remark~\ref{thm:rh_conditional}) is
therefore vacuous: its antecedent cannot be established, because the
criterion it assumes does not hold.
\textbf{Unconditional contribution:}
\textbf{Theorem~\ref{thm:main}: certified strict log-concavity of~$\Phi$ on~$[0,\infty)$
(the theorem is Csordas's~\cite[Theorem~4.5]{csordas2015}; the certification
is independent).}

\medskip\noindent
Thus the source audit does not merely remove an RH corollary; it identifies
the exact missing theorem bridge and then closes the question: no
replacement bridge of the same form can exist, because
Hypothetical Criterion~\ref{hcrit:polya_star} is refuted.
The certified log-concavity result stands independently.
}}
\end{center}

\subsection{Extended source analysis}\label{sec:extended_source}

The following primary sources have been examined or identified as candidates.

\begin{itemize}
\item \textbf{Brand\'{e}n--Chass\'{e} (2016)}~\cite{brandenchasse2016}.
  Section~5 of arXiv:1402.2795v2 explicitly names the \emph{de~Bruijn--Ilieff theorem}:
  ``N.~G.\ de~Bruijn and L.~Ilieff independently proved that if $\phi(t)$ is
  an even entire function and $\phi'(t)\in\mathcal{LP}$, then the Fourier transform of
  $e^{\phi(it)}$ will have only real zeros.''
  The condition is $\phi'(t)\in\mathcal{LP}$ (LP-class condition on the derivative),
  categorically different from log-concavity of the kernel.
\item \textbf{Ilieff (1955)}~\cite{ilieff1955}:
  Brand\'{e}n--Chass\'{e} confirm Ilieff independently proved the same result as
  de~Bruijn Theorem~1 (LP-class condition on $\phi'(t)$).
\item \textbf{De Bruijn (1950) Theorem~20}~\cite{debruijn1950}:
  Brand\'{e}n--Chass\'{e} state they ``prove a
  theorem of de~Bruijn'' [= Theorem~20, by new elementary methods] and separately
  ``extend a theorem of de~Bruijn and Ilieff'' [= Theorem~1], indicating Theorem~20
  is a distinct result from Theorem~1.
  Theorem~20 appears at printed p.~213 of the paper (Duke Math.\ J.\ \textbf{17}
  (1950), 197--226); it is not the paper's final theorem (Theorem~28 follows
  at p.~224), and the text is not paywalled: the Version of Record is freely
  served by the TU/e repository.
  Theorem~20 states: if all roots of the derivative of the polynomial
  $f(t) = \sum_1^N q_i t^i$ ($N$ positive and even) are purely imaginary and
  $q_N > 0$, then $\Psi(z) = \int e^{-f(t)} e^{izt}\,dt$ has real roots only.
  It is a polynomial-exponent theorem; it has no log-concavity hypothesis
  and is not the criterion.
\item \textbf{P\'{o}lya (1930), Quart.\ J.\  Math.\ Oxford 1, 21--34}:
  Contains three conjectures (``Fourier--P\'{o}lya conjecture'') about counting
  nonreal zeros via the behavior of derivatives on the real axis.
  Proved by Ki--Kim (2000).  Entirely unrelated to the log-concavity criterion.
\item \textbf{Coffey--Csordas (2013)}~\cite{csordas2013}:
  zbMATH review (Zbl~1278.30028) explicitly flags as an \textbf{Open Problem}:
  whether a log-concave admissible kernel's moment sequence satisfies the
  Tur\'{a}n inequalities (a necessary condition for only real zeros of the Fourier
  transform).
\item \textbf{Csordas (2015)}~\cite{csordas2015}, Section~3, Theorems~3.5--3.7:
  Proves a new necessary and sufficient condition for the Fourier transform of a
  strictly log-concave admissible kernel to have only real zeros.
  The N\&S condition requires the generalized Laguerre inequalities beyond
  log-concavity alone.
  Explicitly states: ``today, there are no known explicit necessary and sufficient
  conditions that even a `nice' kernel $K(t)$ must satisfy in order that its Fourier
  transform have only real zeros.''
\item \textbf{Csordas--Varga (1989) Theorem~2.2}~\cite{csordas1989}:
  Direct reading reveals this is a restatement of P\'{o}lya Satz~II:
  the hypothesis is that the Mellin transform
  $H(z)=\int_0^\infty t^{z-1}K(t)\,dt$ has only real negative zeros, not
  log-concavity of~$K$.
  Csordas--Varga cite this as ``[23, p.~7]'' where [23] is P\'{o}lya (1927).
  Their Theorem~2.3 is de~Bruijn (1950) Theorem~1, using an LP-class
  condition on $h'(t)=-(\log K)'(t)$, not $(\log K)''\leq 0$.
\end{itemize}

%% =========================================================================
\section{Hypothetical P\'{o}lya-Type Criterion}
\label{sec:polya}

The following criterion is the exact theorem bridge that would be needed to derive RH
from Theorem~\ref{thm:main}.
We have not matched this statement to P\'{o}lya (1927), Csordas--Varga (1989),
de~Bruijn (1950), Ilieff (1955), or another peer-reviewed primary source;
and, as Proposition~\ref{prop:counterexample} shows, no such source can
exist, because the criterion is false.

\begin{hypothetical}[P\'{o}lya-type real-rootedness criterion; false]%
\label{hcrit:polya_star}
Let $K : \mathbb{R} \to \mathbb{R}$ be an even function satisfying:
\begin{enumerate}[label=(\roman*)]
\item $K(t) > 0$ for all~$t$;
\item $K \in L^1(\mathbb{R})$;
\item $(\log K)''(t) \leq 0$ for all~$t \geq 0$;
\item $K(t) = O\bigl(e^{-|t|^{2+\delta}}\bigr)$ for some~$\delta > 0$;
\item $K$ is real analytic on a neighborhood of the origin.
\end{enumerate}
Then the entire function $F(z) = \int_{-\infty}^{\infty} K(t)\,e^{izt}\,dt$
has only real zeros.

\medskip\noindent\emph{Status: FALSE.}
Direct reading of P\'{o}lya (1927) Satz~I and Satz~II gives a preservation theorem
and a Mellin--Fourier theorem, respectively; neither states this criterion.
Csordas--Varga (1989), Theorem~2.2, restates the Mellin--Fourier theorem.
Peer-reviewed literature (Csordas 2015, Theorems~3.5--3.7) shows that within the
admissible-kernel framework, strict log-concavity is not by itself the sufficient
condition for real-rootedness; additional generalized Laguerre inequalities are needed.
Proposition~\ref{prop:counterexample} proves the criterion false as stated:
a kernel satisfying all five hypotheses whose Fourier transform has
non-real zeros.
\end{hypothetical}

\begin{proposition}[The criterion is false]\label{prop:counterexample}
Hypothetical Criterion~\ref{hcrit:polya_star} is false.
The kernel
\[
K(t) = (5 + t^4)\,e^{-t^2 - t^4/1000}
\]
satisfies all five of its hypotheses, but its Fourier transform
$F(z) = \int_{-\infty}^{\infty} K(t) e^{izt}\,dt$ has non-real zeros.
\end{proposition}

\begin{proof}
\emph{Hypotheses.}
$K$ is even and positive, and entire, hence real analytic.
Since $e^{-t^4/1000}$ decays faster than $e^{-|t|^3}$, condition~(iv)
holds with $\delta = 1$.
A direct computation gives
\[
(\log K)''(t)
= \frac{-\,5 - 20\,(t^2 - \tfrac32)^2 - 4t^6 - 2t^8}{(5 + t^4)^2}
  - \frac{3t^2}{250} < 0
\qquad\text{for all } t,
\]
so $K$ is strictly log-concave: condition~(iii) holds.

\emph{Zeros.}
$F$ is entire of order at most $4/3 < 2$ (the quartic damping), hence of
genus at most~$1$; $F$ is even, so its Hadamard factorization has no
linear exponential factor, and $F(0) = m_0 > 0$, where
$m_{2j} := \int_{-\infty}^{\infty} t^{2j} K(t)\,dt$.
If all zeros of $F$ were real, the factorization
$F(z) = m_0 \prod_k (1 - z^2/\alpha_k^2)$ and the expansion
$F(z) = \sum_j (-1)^j m_{2j} z^{2j} / (2j)!$ would give, with
$s_1 = \sum_k \alpha_k^{-2}$ and $s_2 = \sum_k \alpha_k^{-4} \geq 0$,
\[
\frac{m_2}{2 m_0} = s_1,
\qquad
\frac{m_4}{24 m_0} = \frac{s_1^2 - s_2}{2} \leq \frac{s_1^2}{2},
\qquad\text{hence}\qquad
m_0 m_4 \leq 3 m_2^2,
\]
a necessary condition for real-rootedness.
The moments violate it.
With $M_{2j} := \int t^{2j} (5 + t^4) e^{-t^2}\,dt$ the undamped moments
$M_0 = \tfrac{23}{4}\sqrt{\pi}$, $M_2 = \tfrac{35}{8}\sqrt{\pi}$,
$M_4 = \tfrac{165}{16}\sqrt{\pi}$, $M_8 = \tfrac{12495}{64}\sqrt{\pi}$:
the damping factor $e^{-t^4/1000} \leq 1$ gives
$m_2 \leq M_2$, and $e^{-t^4/1000} \geq 1 - t^4/1000$ gives
\[
m_0 \geq M_0 - \tfrac{1}{1000} M_4 = \tfrac{18367}{3200}\sqrt{\pi},
\qquad
m_4 \geq M_4 - \tfrac{1}{1000} M_8 = \tfrac{129501}{12800}\sqrt{\pi}.
\]
Therefore
\[
m_0 m_4 - 3 m_2^2
\;\geq\;
\pi\left(\frac{18367}{3200}\cdot\frac{129501}{12800}
- 3\left(\frac{35}{8}\right)^2\right)
= \frac{26544867}{40960000}\,\pi \;>\; 0,
\]
contradicting the necessary condition.
Hence $F$ has a non-real zero.
(Numerically, a zero is located at $z \approx 2.8086 + 1.3359\,i$,
certified by a winding-number count; and $F(iy) > 0$ for real~$y$,
since the integrand $K(t)e^{-yt}$ is positive, so the non-real zeros
are genuinely off-axis.
Csordas~\cite[Example~3.12]{csordas2015} exhibits the same phenomenon:
log-concavity alone does not force real-rootedness.)
\end{proof}

\begin{remark}[Necessity of analyticity]
Condition~(v) is essential.
The even function $K(t) = e^{-|t|^3}$ satisfies conditions~(i)--(iv) but
fails analyticity at the origin, and its cosine transform has 4~complex zeros
in $[5,15]\times[-5,5]$~\cite[Ex.~2.1]{csordas1989}.
Our computation independently confirms this (Attack~1 in Section~\ref{sec:falsify}).
By contrast, $e^{-t^p}$ for \emph{even integer}~$p$ is real analytic and its cosine
transform has only real zeros.
\end{remark}

\begin{remark}[Hypothesis verification for Hypothetical Criterion~\ref{hcrit:polya_star}]%
\label{rem:polya_hyp}
All five conditions of Hypothetical Criterion~\ref{hcrit:polya_star} are satisfied by~$\Phi$:

\medskip
\noindent\begin{tabular}{@{}llll@{}}
\toprule
Cond.\ of Hyp.\ Crit.~\ref{hcrit:polya_star} & Verified as & Where proved & Reference \\
\midrule
(i) $K(t) > 0$ & $\Phi(u) > 0$ & Prop.~\ref{prop:phi_properties}(i) &
  \cite[Thm.~A]{csordas1989} \\
(ii) $K \in L^1(\mathbb{R})$ & $\Phi \in L^1(\mathbb{R})$ &
  Prop.~\ref{prop:phi_properties}(iii) & \cite[\S2.10]{titchmarsh1986} \\
(iii) $(\log K)'' \leq 0$ & $(\log\Phi)'' < 0$ (strict) & Theorem~\ref{thm:main} &
  \S\S\ref{sec:ia}--\ref{sec:tailbound} \\
(iv) $O(e^{-|t|^{2+\delta}})$ & $\Phi = O(e^{-(\pi-\eta) e^{2u}}) = O(e^{-u^3})$ &
  Prop.~\ref{prop:phi_properties}(iv) & any $0<\eta<\pi$; then $e^{2u}$ outgrows $u^3$ \\
(v) Real analytic near~$0$ & $\Phi$ real analytic on $\mathbb{R}$ &
  Prop.~\ref{prop:phi_properties}(v) & Weierstrass $M$-test \\
(evenness) $K(-t)=K(t)$ & $\Phi(-u) = \Phi(u)$ & Prop.~\ref{prop:phi_properties}(ii) &
  \cite[\S2.10]{titchmarsh1986} \\
\bottomrule
\end{tabular}

\medskip
\noindent See Appendix~\ref{sec:polya_audit} for a hypothesis comparison table.
\end{remark}

%% =========================================================================
\section{Conditional Branches}\label{sec:branches}

The RH implication depends on which form of the P\'{o}lya-type criterion holds.

\paragraph{Branch~A: Strong log-concavity criterion (Hypothetical Criterion~\ref{hcrit:polya_star}).}
Under this criterion, all five conditions verified for~$\Phi$ (Remark~\ref{rem:polya_hyp})
would yield RH.
\emph{This branch is vacuous: the criterion is false
(Proposition~\ref{prop:counterexample}), so the implication has a false
antecedent and yields nothing.
Source audit: Verdict D --- source mismatch.}

\paragraph{Branch~B (Satz~I reading): Preservation theorem.}
Multiplying by an $\mathcal{LP}$-factor preserves real-rootedness but cannot
establish it from scratch.
\emph{Branch~B does not yield RH from log-concavity of~$\Phi$.}

\paragraph{Branch~C (Satz~II reading): Mellin--Fourier theorem.}
The hypothesis is that $H(z) = \int_0^\infty u^{z-1}\Phi(u)\,du$ has only real
negative zeros, which is a separate non-trivial analytic condition not established here.
\emph{Branch~C does not yield RH from log-concavity of~$\Phi$ without additional work.}

\begin{remark}[Conditional RH consequence under Hypothetical Criterion~\ref{hcrit:polya_star}]%
\label{thm:rh_conditional}
\textbf{This remark is not an unconditional mathematical consequence of the results
proved in this paper.  It is vacuous.}

If Hypothetical Criterion~\ref{hcrit:polya_star} held, then the
conditions verified for~$\Phi$ in Remark~\ref{rem:polya_hyp} would yield: the Fourier
transform $F(z) = 2\Xi(z)$ has only real zeros, so~$\Xi$ has only real zeros, which
is equivalent to RH~\cite{titchmarsh1986}.
But the criterion does not hold: Proposition~\ref{prop:counterexample}
proves it false.
The implication is therefore vacuous, no route to RH exists through it,
and \textbf{no proof of RH is claimed.}
The remark is retained, relabeled, so that the earlier conditional
framing in the literature trail remains traceable.
\end{remark}

\paragraph{The audit's question, answered.}
The source audit reduced the possible RH implication of Theorem~\ref{thm:main}
to the question whether Hypothetical Criterion~\ref{hcrit:polya_star} is true.
That question is now answered in the negative
(Proposition~\ref{prop:counterexample}).
Any future bridge from log-concavity of~$\Phi$ to real-rootedness
of~$\Xi$ must use additional hypotheses (for instance the generalized
Laguerre inequalities of Csordas~\cite{csordas2015}); no criterion with
exactly the five hypotheses above can apply, to~$\Phi$ or to any kernel.

%% =========================================================================
\section{Proof-Critical Dependencies}\label{sec:deps}

\begin{center}\footnotesize
\begin{tabular}{@{}p{4.0cm}p{2.2cm}p{3.8cm}l@{}}
\toprule
Claim & Type & Primary source & Status \\
\midrule
$\varphi_1(u) > 0$ for $u \geq 0$ & Lean & \texttt{phi1\_pos} & PROVED \\
$(\log\varphi_1)''(u) < 0$ for $u \geq 0$ & Lean & \texttt{log\_phi1\_d2\_neg} & PROVED \\
$\Phi(u) > 0$ for $u \geq 0$ & Cited & \cite[Thm.~A]{csordas1989} & CITED \\
$\Phi(-u) = \Phi(u)$ & Cited & \cite[\S2.10]{titchmarsh1986} & CITED \\
$\Phi \in L^1(\mathbb{R})$ & Cited & \cite[\S2.10]{titchmarsh1986} & CITED \\
$\Phi$ real analytic on $\mathbb{R}$ & Structural & Weierstrass $M$-test & PROVED \\
$\Phi = O(e^{-(\pi-\eta) e^{2u}})$, $0<\eta<\pi$ & Structural & \texttt{phi1\_decay\_bound} & PROVED \\
$Q_\Phi(u) < 0$ on $[0,1]$ & IA cert & Arb/FLINT (primary), mpmath.iv (cross-check) & CERTIFIED \\
$(\log\Phi)'' < 0$ on $[1,3]$ & IA cert & Arb/FLINT (primary), mpmath.iv (cross-check) & CERTIFIED \\
$(\log\Phi)'' < 0$ for $u \geq 3$ & Monotonicity & $\varepsilon^*(3)<10^{-1636}$ & PROVED \\
Hyp.\ Crit.~\ref{hcrit:polya_star} & None & No peer-reviewed source &
  \textbf{REFUTED} \\
RH $\Leftrightarrow$ $\Xi$ real zeros & Cited & \cite[\S2.1]{titchmarsh1986} & CITED \\
\bottomrule
\end{tabular}
\end{center}
\medskip

\noindent\textbf{Status key:}
PROVED = formal (Lean) or elementary argument;
CERTIFIED = rigorous computational certificate;
CITED = external theorem with explicit hypothesis verification;
REFUTED = disproved by the counterexample of
Proposition~\ref{prop:counterexample}.

\medskip
\noindent\textbf{Only the PROVED, CERTIFIED, and ordinary CITED entries are used in the
unconditional proof of Theorem~\ref{thm:main}.
The REFUTED item is not used in any unconditional theorem.}
The REFUTED entry is not a defect in the proof of Theorem~\ref{thm:main};
it marks the boundary between the unconditional log-concavity theorem proved here
and the RH interpretation that the refuted criterion would have supported.

%% =========================================================================
\section{Falsification Testing}
\label{sec:falsify}

We subjected every proof-chain link to 36~systematic tests in 7~batches.
No test detected any inconsistency.

\begin{center}
\small
\begin{tabular}{clp{5.8cm}l}
\toprule
\# & Target & Attack & Result \\
\midrule
1 & Analyticity & $e^{-t^3}$ cosine transform complex zeros & 4 found \\
2 & $\Phi > 0$ & Search 10{,}001 points on $[0,1]$ & $\min = 5.5 \times 10^{-7}$ \\
4 & $Q_\Phi < 0$ & 7{,}001 points (dense + random + adversarial) & All negative \\
6 & Convention & $\int\Phi\,du = \xi(1/2)$ & Ratio $= 1.000$ \\
7 & Perturbation diagnostic & Compute legacy $C$ at $u = 1$ & $C = 204$ \\
10 & IA correctness & Containment: $\pi$, $e$, $e^{-\pi}$, compound & All pass \\
22 & Product rule & Full $\varphi_1''$ vs.\ \texttt{mpmath.diff} & $< 10^{-15}$ \\
24 & $\Xi(5)$ & $\int\Phi\cos(5u)\,du$ vs.\ $\xi(\tfrac{1}{2}+5i)$ & Ratio $= 1.000$ \\
27 & P\'{o}lya on $e^{-\cosh t}$ & Argument-principle count & Only real \\
31 & Adversarial $Q$ & 15{,}001 points on $[0.5, 2.0]$ & All negative \\
\bottomrule
\end{tabular}
\end{center}
\noindent Attack~12 historically detected a real bug (wrong $g''$ coefficient).
The remaining 24~attacks cover: $\Phi$ evenness, decay, smoothness, circularity,
series convergence, end-to-end checks at $\gamma_2$, and extended attacks 33--36
($10^5$-point dense $Q_\Phi$ scan; P\'{o}lya conditions (i)--(v); Jensen
polynomial hyperbolicity; $\Lambda = 0$ consistency).

%% =========================================================================
\section{Reproducibility}
\label{sec:reproduce}

\paragraph{Interval arithmetic rigor.}
The primary rigorous interval certificate for~$[0,1]$ is the Arb/FLINT
ball-arithmetic verification, implemented via \texttt{python-flint}.
Arb uses ball arithmetic with rigorous error bounds and directed enclosures for
elementary functions.
The Arb/FLINT ball-arithmetic implementation is the primary rigorous interval
certificate for both $[0,1]$ and $[1,3]$.
The \texttt{mpmath.iv} implementation is retained as an independent cross-check
for both interval certifications.
Constants $\pi$ and $e$ are enclosed by rigorous interval or ball enclosures
before entering interval operations.
Elementary functions use directed outward-rounding enclosures.
Subinterval endpoints are represented as exact multi-precision floats.
\textbf{The interval arithmetic certification does not rely on point sampling:
each subinterval is evaluated as an interval object or ball enclosure and
certified by a negative upper bound.}

\begin{center}\footnotesize
\begin{tabular}{@{}p{3.5cm}p{1.6cm}p{4.5cm}@{}}
\toprule
Purpose & Verifier & Script and command \\
\midrule
Certify $Q_\Phi<0$ on $[0,1]$ &
  Arb/FLINT (primary) &
  \texttt{proof/verify\_logconcavity\_arb.py} \\
  & & \texttt{python proof/verify\_logconcavity\_arb.py} \\[3pt]
Cross-check $Q_\Phi<0$ on $[0,1]$ &
  mpmath.iv &
  \texttt{proof/verify\_logconcavity\_rigorous.py} \\
  & & \texttt{python proof/verify\_logconcavity\_rigorous.py} \\[3pt]
Certify $(\log\Phi)''<0$ on $[1,3]$ &
  Arb/FLINT (primary) &
  \texttt{proof/verify\_ia\_1\_to\_3\_arb.py} \\
  & & \texttt{python proof/verify\_ia\_1\_to\_3\_arb.py} \\[3pt]
Cross-check $(\log\Phi)''<0$ on $[1,3]$ &
  mpmath.iv &
  \texttt{proof/verify\_ia\_1\_to\_3.py} \\
  & & \texttt{python proof/verify\_ia\_1\_to\_3.py} \\[3pt]
Full verification pipeline &
  both &
  \texttt{verify.py} (7 steps, fail-closed) \\
  & & \texttt{python verify.py} \\[3pt]
Falsification tests &
  both &
  \texttt{falsify.py} \quad \texttt{python falsify.py} \\
\bottomrule
\end{tabular}
\end{center}

All commands were run from commit \texttt{8a670d7} on 2026-11-22.

\paragraph{Lean~4 dependency audit.}
The Lean formalization is a dependency audit and algebraic-core verification.
It proves the dominant-term positivity and curvature statements and records the
interval arithmetic, Fourier representation, classical $\Phi$ properties, and
the hypothetical P\'{o}lya-type criterion as explicit axioms.
13~project-specific axioms; 6~proved auxiliary theorems, plus the
conditional assembly theorem \texttt{riemann\_hypothesis} proved from
those axioms; 0~\texttt{sorry}.
Since the criterion recorded by the axiom \texttt{polya\_theorem} is
false (Proposition~\ref{prop:counterexample}), the assembly theorem is
vacuous: it certifies the dependency structure, not RH.
Machine-checked against Mathlib commit \texttt{a6276f4c} (0~errors, 0~correctness warnings).

\begin{center}\footnotesize
\begin{tabular}{@{}lp{3.0cm}p{2.5cm}p{3.5cm}@{}}
\toprule
Statement & Lean status & Role & Note \\
\midrule
\texttt{phi1\_pos} & proved & algebraic core & unconditional \\
\texttt{log\_phi1\_d2\_neg} & proved & algebraic core & unconditional \\
\texttt{ia\_verification\_0\_to\_1} & axiom & computational cert & certifies $Q_\Phi < 0$ on $[0,1]$ \\
\texttt{ia\_verification\_1\_to\_3} & axiom & computational cert & certifies $(\log\Phi)'' < 0$ on $[1,3]$ \\
\texttt{polya\_theorem} & axiom; criterion REFUTED & hypothetical criterion &
  assembly vacuous; not used unconditionally \\
\texttt{rh\_iff\_xi\_real} & axiom / cited & conditional only &
  \cite[\S2.1]{titchmarsh1986} \\
\bottomrule
\end{tabular}
\end{center}

%% =========================================================================
\section{Discussion}
\label{sec:discussion}

\paragraph{Strengths.}
The algebraic core is machine-checked in Lean~4.
The $[0,1]$ certification is reproduced independently using Arb/FLINT ball arithmetic.
The mpmath.iv certificate uses 52{,}898 subintervals; the Arb/FLINT reproduction
uses its own interval partition and precision settings, reported in
\texttt{results/verify\_logconcavity\_arb.json}.
The extended $[1,3]$ certification avoids catastrophic cancellation by
working with $(\log\Phi)''$ directly, using 1{,}813 adaptive overlapping
balls, with Arb/FLINT as primary certificate and
mpmath.iv as cross-check.

\paragraph{Relation to prior claims.}
A preprint by Gershon (April 2026)~\cite{gershon2026} establishes the same
log-concavity result but contains a presentation error in the perturbation bound.
Our treatment replaces the presentation-level perturbation constant with an
interval-certified $B(u)$ bound on~$[1,3]$ and an explicit analytic far-tail bound
for $u\ge3$, thereby extending the certification to the full half-line.
Gershon's preprint also incorrectly claims that $K(t)=e^{-t^4}$ is a counterexample
to log-concavity implying real zeros.
This is wrong: Csordas--Varga~\cite[Ex.~2.1]{csordas1989} explicitly state that
$\int e^{-t^p}\cos(zt)\,dt$ has only real zeros for every \emph{even integer} $p$,
confirmed independently by our computation (Attack~27).
The kernel $K(t)=e^{-|t|^3}$ is a counterexample only to the criterion
without its analyticity condition~(v), which it fails at the origin.
The sharp counterexample, satisfying all five conditions including~(v),
is $K(t) = (5+t^4)e^{-t^2-t^4/1000}$ (Proposition~\ref{prop:counterexample}).

Contextual implications are in Appendix~\ref{sec:context}.

%% =========================================================================
\section{Limitations}
\label{sec:limitations}

\begin{enumerate}[leftmargin=2em,itemsep=2pt]
\item The log-concavity-to-real-rootedness criterion is false
  (Proposition~\ref{prop:counterexample}), and direct source audit gives
  a mismatch with P\'{o}lya Satz~I/II and Csordas--Varga Theorem~2.2.
  The result proved here is log-concavity of~$\Phi$, not RH.
\item The interval arithmetic certificates are not formalized in Lean.
  Axioms \texttt{ia\_verification\_0\_to\_1} and \texttt{ia\_verification\_1\_to\_3}
  remain axioms, not derivations within the proof assistant.
  They are external certificates, with Arb/FLINT primary and mpmath.iv cross-checks.
\item The proof of classical properties of~$\Phi$ relies on standard references.
\item The work has not undergone specialist peer review.
\item The auxiliary Hilfss\"{a}tze of P\'{o}lya (1927) have not all been
  fully transcribed.
  The one that has been cited under the name ``P\'{o}lya's Hilfssatz~II''
  is identified and closed (Cardon~\cite{cardon2005}, Lemma~2, sourced to
  P\'{o}lya (1926)): a shift-sum preservation lemma, not the criterion.
  The remaining untranscribed Hilfss\"{a}tze cannot rehabilitate the
  criterion in any case, since the criterion is false.
\end{enumerate}

%% =========================================================================
\bibliographystyle{plain}
\bibliography{references}

%% =========================================================================
\appendix

%% =========================================================================
\section{Contextual Results}\label{sec:context}

The following results are consequences of or evidence consistent with RH but
\emph{do not form part of the proof chain} in Section~\ref{sec:deps}.

\paragraph{(i) Jensen polynomial hyperbolicity.}
\cite{griffin2019} establish that $J^d_n(X)$ is hyperbolic for each fixed~$d$.
If Remark~\ref{thm:rh_conditional} were established unconditionally, this would
extend to all $d, n \geq 0$ simultaneously.
Numerically verified for $k = 1,\ldots,8$ ($d=2$) and $n = 1,\ldots,7$ ($d=3$).

\paragraph{(ii) De Bruijn--Newman constant.}
Conditional on the RH consequence, $\Lambda = 0$ (combining Rodgers--Tao~\cite{rodgers2020}
and Platt--Trudgian).

\paragraph{(iii) CvS spectral framework.}
Eigenvalue plateau at $N = 10$: $\log_{10}|\lambda_{\min}(c)|$ reaches $-34.2$ at $c = 53$.
Included as independent numerical evidence only.

\paragraph{(iv) AEE assessment.}
The AEE framework~\cite{specsmith} assigns 0.169/1.000 (non-critical failures: Lean
certificates not derived within Lean, no peer review, Verdict~D source mismatch).

%% =========================================================================
\section{Proof Certificate Table}\label{sec:cert}

\begin{center}\footnotesize
\begin{tabular}{@{}lp{2.8cm}p{2.4cm}lp{2.8cm}@{}}
\toprule
Step & Claim & Type & Status & Script \\
\midrule
C1 & $\varphi_1 > 0$ & Lean & PROVED & \texttt{phi1\_pos} \\
C2 & $(\log\varphi_1)'' < 0$ & Lean & PROVED & \texttt{log\_phi1\_d2\_neg} \\
C3 & $\Phi > 0$, even, $\in L^1$ & Classical & CITED & \cite{titchmarsh1986,csordas1989} \\
C4 & $Q_\Phi < 0$ on $[0,1]$ & IA (Arb/FLINT) & CERT & \texttt{verify\_logconcavity\_arb.py} \\
C4b & $Q_\Phi < 0$ on $[0,1]$ & IA (mpmath.iv) & CERT & \texttt{verify\_logconcavity\_rigorous.py} \\
C5 & $(\log\Phi)'' < 0$ on $[1,3]$ & IA (Arb/FLINT) & CERT & \texttt{verify\_ia\_1\_to\_3\_arb.py} \\
C5b & $(\log\Phi)'' < 0$ on $[1,3]$ & IA (mpmath.iv) & CERT & \texttt{verify\_ia\_1\_to\_3.py} \\
C6 & $(\log\Phi)'' < 0$, $u\geq 3$ & Monotonicity & PROVED & \texttt{verify\_weighted\_comparisons.py} \\
C7 & $\Phi$ real analytic & Structural & PROVED & Weierstrass $M$-test \\
C8 & $\Phi = O(e^{-(\pi-\eta) e^{2u}})$ & Structural & PROVED & \texttt{phi1\_decay\_bound} \\
C9 & Hyp.\ Crit.~\ref{hcrit:polya_star} & None &
  \textbf{REFUTED} & Prop.~\ref{prop:counterexample} \\
C10 & RH $\Leftrightarrow$ $\Xi$ real zeros & Cited & CITED & \cite[\S2.1]{titchmarsh1986} \\
\bottomrule
\end{tabular}
\end{center}
\noindent\footnotesize CERT = CERTIFIED (rigorous computational certificate);
C4/C5 are primary, C4b/C5b are cross-checks.\normalsize

\noindent\textbf{Only C1--C8 and C10 are used in the unconditional proof of
Theorem~\ref{thm:main}.
Item~C9 (REFUTED) is not used in any unconditional theorem.}

\paragraph{Proof-critical fields in \texttt{results/verify\_ia\_1\_to\_3\_arb.json}.}
The following fields are directly cited in Appendix~\ref{sec:appE}:
\begin{center}\small
\begin{verbatim}
{
  "claim": "(log Phi)''(u) < 0 on [1,3]",
  "method": "Arb/FLINT ball arithmetic via python-flint; direct DeltaQ
             enclosure (finite sum n=1..10) plus explicit derivative-aware
             tail majorants for n>10; sign test (W_1 + B) < 0 in Arb",
  "coverage": "[0.99,3.01]",
  "intervals_checked": 1813,
  "precision_bits": 256,
  "finite_terms": 10,
  "tail_allowance_B": "6.55494e-1174",
  "tail_bound_n_gt_10": "|R''| tail < 6.55494e-1174; eps* tail < 1.78467e-2515",
  "max_B_upper": "8.538902e-24",
  "minimum_margin": "9.129797e+01",
  "passed": true
}
\end{verbatim}
\end{center}
\noindent The previous \texttt{mpmath.iv} output \texttt{results/verify\_ia\_1\_to\_3.json}
is retained as an independent cross-check.

\paragraph{Far-tail bound ($u\geq3$).}
The far-tail domination $|W_{\mathrm{tail}}(u)|\leq2600\,e^{4u}\varepsilon^*(u)$ for $u\geq3$
is proved analytically in Appendix~\ref{sec:appE7} (Lemmas~\ref{lem:weighted_comparisons}--\ref{lem:fartail_mono})
and does not depend on a JSON certificate.\\[2pt]
Summary of far-tail analytic bound:
\begin{center}\small
\begin{verbatim}
{
  "far_tail_claim": "|Wtail(u)| <= 2600 * e^{4u} * eps_star(u) for u >= 3",
  "C_tail": 2600,
  "eps_star_3_upper": "1.655889e-1650",
  "far_tail_margin_at_3": "W1(3) < -1000, C_tail*e^12*eps_star(3) < 1e-1629",
  "analytic_proof": "Appendix E.7",
  "passed": true
}
\end{verbatim}
\end{center}

\medskip\noindent
Full SHA256 hashes, commit \texttt{8a670d7}, date 2026-11-22:

\begin{verbatim}
results/verify_logconcavity_rigorous.json
CE97B742356F14AFFD5A1EF93DE3EDD7B63D52CE7924F37BFF4DA31835FCF2F8

results/verify_logconcavity_arb.json
646698D243299DF7ADB0542F7378EBD8EE30F299D70600E8EF90093382BC6557

results/verify_ia_1_to_3.json
53714517AD033DC65636B1E30F2866C9D0631FD232B496FD78DCE04D9317C95D

results/verify_ia_1_to_3_arb.json
BB889E9A1A0FD0CE75F5769FD7BE1DC34B1AF7C7987B0415779ABF8959C5A592

verification/proof_certificate_v2.json
8B538345D589638A7AB9534470085CF531DACABD57D9143A2BF327C5CB2DC192
\end{verbatim}

%% =========================================================================
\section{P\'{o}lya Source Analysis}\label{sec:polya_audit}

The source audit is in Section~\ref{sec:polya_audit_main}.
For reference, the hypothesis comparison between P\'{o}lya (1927) Satz~II and
Csordas--Varga (1989) Theorem~2.2:

\begin{center}\small
\begin{tabular}{@{}p{4.2cm}p{4.8cm}p{2.6cm}@{}}
\toprule
P\'{o}lya 1927 Satz~II & Csordas--Varga 1989 Thm~2.2 & Our $\Phi$ satisfies \\
\midrule
$f$ for $t>0$, real-valued, abs.\ integrable & Even $K > 0$, in $L^1$ & Props~(i)--(iii) \\
$|f(t)|<Be^{-t^{1+\beta}}$ & $K = O(e^{-|t|^{2+\delta}})$, $\delta>0$ & Prop.~(iv) \\
$f$ analytic near $t=0$ & $K$ real analytic near origin & Prop.~(v) \\
\textbf{$H(z)$ has only real negative zeros} &
  \textbf{$H(z)$ has only real negative zeros} &
  Not established for~$\Phi$ \\
Conc: $\int f(t^{2q})e^{izt}dt$ has real zeros &
  Conc: $\int K(t)e^{izt}dt$ has real zeros & --- \\
\bottomrule
\end{tabular}
\end{center}

The citation chain is broken at both ends: P\'{o}lya (1927) Satz~I/II do not state
the log-concavity criterion, and Csordas--Varga (1989) Theorem~2.2 is P\'{o}lya Satz~II
(Mellin--Fourier theorem).

%% =========================================================================
\section{Verification Checklist}\label{sec:checklist}

\begin{itemize}[leftmargin=2em,itemsep=2pt]
\item[$\square$] The paper never states or implies that RH is proved.
\item[$\square$] The log-concavity theorem (Theorem~\ref{thm:main}) is clearly
  identified as the main unconditional result.
\item[$\square$] The P\'{o}lya-type criterion is labeled hypothetical and refuted
  (Hypothetical Criterion~\ref{hcrit:polya_star},
  Proposition~\ref{prop:counterexample}), not as an established theorem.
\item[$\square$] Satz~I and Satz~II are transcribed and translated
  (Section~\ref{sec:polya_audit_main}).
\item[$\square$] The Hilfssatz~II citation trail is closed
  (Cardon~\cite{cardon2005}, Lemma~2; P\'{o}lya (1926)); the remaining
  uninspected Hilfss\"{a}tze are explicitly marked as uninspected
  (Section~\ref{sec:verdict}).
\item[$\square$] The $[1,3]$ perturbation proof is complete in Appendix~\ref{sec:appE}.
\item[$\square$] All certificate hashes are full SHA256 hashes (Appendix~\ref{sec:cert}).
\item[$\square$] All scripts, outputs, and Lean axiom names match the repository
  (Section~\ref{sec:reproduce}).
\item[$\square$] Lean claims are limited to algebraic-core verification and dependency
  tracking (Section~\ref{sec:reproduce}).
\item[$\square$] RH appears only in a conditional remark (Remark~\ref{thm:rh_conditional}).
\end{itemize}

%% =========================================================================
\section{Details of the Perturbation and Tail Bounds}
\label{sec:appE}

This appendix provides complete written derivations for all estimates underlying
Lemma~\ref{lem:wtail} and Theorem~\ref{thm:extended}.
Scripts cited here serve only as reproducibility checks and interval-arithmetic
certificates for stated numerical constants; the algebraic argument for every bound
is given below.

\subsection{Notation and scaling}

Write $\Phi(u) = 4\Phi_{\mathrm{red}}(u)$ where
$\Phi_{\mathrm{red}} := \varphi_1 + R$ and $R := \sum_{n\geq 2}\varphi_n$.
Since $(\log(4f))'' = (\log f)''$ for any positive~$f$, the factor~4 does not
affect $(\log\Phi)''$.
All subsequent analysis uses $\Phi_{\mathrm{red}}$.

Define:
\begin{align*}
h(u) &:= 2\pi e^{2u} - 3, \quad h(u) \geq h(0) = 2\pi-3 > 0 \text{ for } u\geq 0,\\
W_1(u) &:= (\log\varphi_1)''(u) = -\tfrac{24\pi e^{2u}}{h(u)^2} - 4\pi e^{2u},\\
W(u) &:= (\log\Phi_{\mathrm{red}})''(u),\\
W_{\mathrm{tail}}(u) &:= W(u) - W_1(u),\\
Q_f(u) &:= f''(u)f(u) - f'(u)^2 \quad\bigl((\log f)'' = Q_f/f^2\bigr),\\
\Delta Q(u) &:= Q_{\Phi_{\mathrm{red}}}(u) - Q_{\varphi_1}(u),\\
\varepsilon^*(u) &:= 2\sum_{n=2}^{\infty} n^4 e^{-\pi(n^2-1)e^{2u}}.
\end{align*}

Theorem~\ref{thm:algebraic} gives $W_1(u) < -4\pi e^{2u}$.
In particular $W_1(1) \approx -93.15$ and $\varepsilon^*(1) \approx 1.82\times 10^{-29}$
(both confirmed by interval arithmetic, \texttt{verify\_ia\_1\_to\_3.py}).

\subsection{Expansion of $\Delta Q$}

\begin{lemma}\label{lem:deltaQ}
$\Delta Q = R''\varphi_1 + \varphi_1''R - 2\varphi_1'R' + R''R - (R')^2$.
\end{lemma}

\begin{proof}
Expand $Q_{\Phi_{\mathrm{red}}} = (\varphi_1''+R'')(\varphi_1+R) - (\varphi_1'+R')^2$:
\[
= \varphi_1''\varphi_1 - (\varphi_1')^2
  + (R''\varphi_1 + \varphi_1''R - 2\varphi_1'R')
  + (R''R - (R')^2)
= Q_{\varphi_1} + \Delta Q.
\]
\end{proof}

\subsection{Explicit derivative formulas on $[1,3]$}

The certification evaluates the exact closed-form derivatives below
directly in ball or interval arithmetic; no derivative-ratio bounds are
used anywhere in the $[1,3]$ certificate.
(An earlier version of this appendix stated uniform bounds
$|\varphi_n'| \leq 4\,n^4 e^{-\pi(n^2-1)e^{2u}}|\varphi_1|$ and
$|\varphi_n''| \leq 8\,n^4 e^{-\pi(n^2-1)e^{2u}}|\varphi_1|$; those
bounds are false --- the true ratios grow like $n^6$ and $n^8$ ---
and the lemma asserting them has been removed.  Nothing in the
certificate depended on them: the code has always evaluated the closed
forms directly.)

\begin{lemma}[Closed-form derivatives]\label{lem:derivbounds}
For every $n \geq 1$ and all $u$:
\[
g_n(u) = 2\pi^2 n^4 e^{9u/2} - 3\pi n^2 e^{5u/2}, \quad
g_n'(u) = 9\pi^2 n^4 e^{9u/2} - \tfrac{15}{2}\pi n^2 e^{5u/2},
\]
\[
g_n''(u) = \tfrac{81}{2}\pi^2 n^4 e^{9u/2} - \tfrac{75}{4}\pi n^2 e^{5u/2}.
\]
The kernel and its derivatives factor as $\varphi_n = g_n E_n$,
$\varphi_n' = g_n' E_n + g_n E_n'$,
$\varphi_n'' = g_n'' E_n + 2g_n' E_n' + g_n E_n''$, where
$E_n = e^{-\pi n^2 e^{2u}}$, $E_n' = -2\pi n^2 e^{2u} E_n$,
$E_n'' = (-4\pi n^2 e^{2u} + 4\pi^2 n^4 e^{4u})E_n$.
Note $E_n = E_1 \cdot e^{-\pi(n^2-1)e^{2u}}$.

This gives the exact relative-derivative expressions:
\[
\frac{\varphi_n'}{\varphi_n} = \frac{g_n'}{g_n} - 2\pi n^2 e^{2u}, \qquad
\frac{\varphi_n''}{\varphi_n} = \frac{g_n''}{g_n}
  - 4\pi n^2 e^{2u}\frac{g_n'}{g_n} - 4\pi n^2 e^{2u} + 4\pi^2 n^4 e^{4u}.
\]
\end{lemma}

\begin{proof}
Each identity is the product rule applied to $\varphi_n = g_n E_n$
with the displayed derivatives of $g_n$ and $E_n$; the relative forms
follow by dividing by $\varphi_n = g_n E_n$.
\end{proof}

For $k=0$ there is a clean ratio bound, used in
Lemma~\ref{lem:Wtail_bound}: $|\varphi_n|/|\varphi_1|
= (|g_n|/|g_1|)\cdot(E_n/E_1)
\leq 2n^4 e^{-\pi(n^2-1)e^{2u}}$ since $|g_n/g_1| \leq 2n^4$
(Proposition~\ref{prop:tail_ratio}).

The infinite range $n\geq 2$ is handled by splitting into a finite
evaluated part and an explicit majorant tail.
The verifier evaluates the explicit expressions for $n=1,\ldots,10$
directly on each ball (\texttt{finite\_terms = 10} in the JSON output).
For $n > 10$ and $u \in [0.99, 3.01]$, the code bounds
$|\varphi_n^{(k)}(u)| \leq C_k\, n^{4+2k} e^{9\cdot 3.01/2}
e^{-\pi n^2 e^{2\cdot 0.99}}$ with explicit constants
$C_0 = 30$, $C_1 = 100 + 60\pi e^{6.02}$,
$C_2 = 460 + 460\pi e^{6.02} + 30\,(4\pi e^{6.02} + 4\pi^2 e^{12.04})$,
obtained from the displayed formulas by the triangle inequality
(growth factors at the right endpoint, decay at the left).
The majorants are summed for $n = 11,\ldots,80$ in ball arithmetic,
and the tail beyond $n = 80$ is at most twice the first neglected
majorant, since the majorant ratio
$((n{+}1)/n)^{4+2k} e^{-\pi(2n+1)e^{1.98}}$ is certified $< 1/2$ at
$n = 80$.
The resulting allowances are
$\tau_0 \leq 6.67\times 10^{-1185}$,
$\tau_1 \leq 2.09\times 10^{-1179}$,
$\tau_2 \leq 6.56\times 10^{-1174}$
for $|\sum_{n>10}\varphi_n^{(k)}|$, $k = 0,1,2$;
the $\varepsilon^*$ tail ($n > 15$) is bounded by
$1.79\times 10^{-2515}$ by the same method.
(An earlier version asserted a $10^{-900}$ tail bound in the JSON
metadata without computing it; the bounds above are the ones the code
actually evaluates.)

\begin{corollary}\label{cor:Rbounds}
For all $u \geq 0$: $|R(u)| \leq \varepsilon^*(u)\,|\varphi_1(u)|$.
\end{corollary}

\begin{proof}
$|R(u)| \leq \sum_{n\geq2}|\varphi_n(u)|
\leq \sum_{n\geq2} 2n^4 e^{-\pi(n^2-1)e^{2u}}|\varphi_1(u)|
= \varepsilon^*(u)|\varphi_1(u)|$ by the $k=0$ ratio bound.
\end{proof}

\noindent\emph{Note.}
Only the $k=0$ bound of Corollary~\ref{cor:Rbounds} is used
(in Lemma~\ref{lem:Wtail_bound}).
The $\Delta Q$ term in that lemma is enclosed directly from its
definition in Lemma~\ref{lem:Bsign}, not through derivative-ratio
estimates for $R'$ or $R''$.

\subsection{Bounding $W_{\mathrm{tail}}$}

\begin{lemma}\label{lem:Wtail_bound}
For $u \geq 0$, with $\rho := R(u)/\varphi_1(u) \leq 1/50$:
\[
|W_{\mathrm{tail}}(u)| \leq B(u) :=
\frac{|\Delta Q(u)|}{\varphi_1(u)^2} + 3\varepsilon^*(u)|W_1(u)|.
\]
\end{lemma}

\begin{proof}
Express $W_{\mathrm{tail}} = Q_{\Phi_{\mathrm{red}}}/\Phi_{\mathrm{red}}^2 - Q_{\varphi_1}/\varphi_1^2$:
\[
W_{\mathrm{tail}}
= \frac{\Delta Q\cdot\varphi_1^2 - Q_{\varphi_1}(2\varphi_1 R+R^2)}
       {\Phi_{\mathrm{red}}^2\varphi_1^2}.
\]
Since $Q_{\varphi_1}<0$, $R>0$: both terms satisfy
$|W_{\mathrm{tail}}| \leq {|\Delta Q|}/{\varphi_1^2} + |W_1|\cdot(2\varphi_1 R+R^2)/\varphi_1^2$,
where we used $\Phi_{\mathrm{red}}^2 \geq \varphi_1^2$.
Since $R \leq \varepsilon^*\varphi_1$ (Corollary~\ref{cor:Rbounds}) and
$R \leq \varphi_1$ (as $\rho \leq 1/50$), we have $R^2 \leq \varphi_1 R$ and
$(2\varphi_1 R+R^2)/\varphi_1^2 \leq 3R/\varphi_1 \leq 3\varepsilon^*$.
\end{proof}

\subsection{Interval-arithmetic certificate on $[1,3]$}

\begin{lemma}[Interval-arithmetic certificate]\label{lem:Bsign}
For every $u\in[1,3]$, the ball-arithmetic upper bound $B(u)$ satisfies:
\[
B(u) \leq 8.54\times 10^{-24},
\]
and the sign certificate holds:
\[
W_1(I) + B(I) < 0
\]
(a comparison of Arb balls) for every ball $I$ in the adaptive cover of
$[0.99, 3.01]$ described below.

This is a finite ball-arithmetic certificate.
The expressions being enclosed are the explicit closed-form formulas
for $\Delta Q$, $\varphi_1$, $\varepsilon^*$, and $W_1$ stated above.
No point sampling is used; each ball is evaluated as a ball object and
certified by a negative Arb comparison.
\end{lemma}

\begin{proof}[Certificate details]
The Arb/FLINT verifier
(\texttt{proof/verify\_ia\_1\_to\_3\_arb.py}, function \texttt{certify\_checkpoint\_arb})
works on an adaptive grid of 1{,}813 balls covering $[0.99, 3.01]$:
the half-width at left endpoint~$u$ is
$\min(0.01,\, 0.35/(2\pi e^{2u}))$, so the exponent $-\pi e^{2u}$ of the
dominant term varies by at most $0.35$ across any one ball and the ball
evaluation of $\varphi_1$ stays positive and tight on the whole range.
On each ball $I$:
\begin{enumerate}[leftmargin=2em,itemsep=1pt]
\item Evaluates $\varphi_1, \varphi_1', \varphi_1''$ and the finite tail
  sums $A = \sum_{n=2}^{10}\varphi_n$ (with derivatives) as Arb balls
  from the closed forms of Lemma~\ref{lem:derivbounds}.
\item Encloses $\Delta Q$ directly from its definition
  (Lemma~\ref{lem:deltaQ}) applied to the finite sums, and adds the
  $n > 10$ tail contribution through the explicit allowances
  $\tau_0, \tau_1, \tau_2$ of the derivative-formulas subsection above:
  $|\Delta Q| \leq |\Delta Q_{\mathrm{fin}}|
  + \tau_2|\varphi_1| + |\varphi_1''|\tau_0 + 2|\varphi_1'|\tau_1
  + \tau_2|A| + |A''|\tau_0 + 2|A'|\tau_1 + \tau_2\tau_0 + \tau_1^2$,
  every factor a rigorous \texttt{abs\_upper} ball bound.
\item Evaluates $\varepsilon^*(I)$ as the finite ball sum $n = 2,\ldots,15$
  plus the explicit tail allowance $1.79\times 10^{-2515}$, and
  $W_1(I)$ from its closed form.
\item Forms the point bound
  $B(I) = |\Delta Q|/\mathrm{lower}(\varphi_1)^2
  + 3\,\mathrm{upper}(\varepsilon^*)\,|W_1|$
  and certifies the Arb ball comparison $W_1(I) + B(I) < 0$.
\end{enumerate}
All 1{,}813 checkpoints pass.
Maximum certified $B$ upper bound: $8.54\times 10^{-24}$
(at the leftmost checkpoint, $u\approx1.0$).
Minimum margin $(-W_1 - B)$: $91.29$ (at $u\approx1.0$).
Output: \texttt{results/verify\_ia\_1\_to\_3\_arb.json},
SHA256:

{\small\verb|BB889E9A1A0FD0CE75F5769FD7BE1DC34B1AF7C7987B0415779ABF8959C5A592|}

\noindent The \texttt{mpmath.iv} verifier
\texttt{proof/verify\_ia\_1\_to\_3.py} is an independent cross-check
(101 uniform checkpoints; it encloses the same $\Delta Q$ with the same
majorant tails in interval arithmetic): maximum $B \leq 1.12\times 10^{-23}$,
minimum margin $91.29$; output \texttt{results/verify\_ia\_1\_to\_3.json}
(SHA256 in Appendix~\ref{sec:cert}).

\noindent\emph{History.}
Earlier versions of this certificate evaluated a surrogate bound
($2500\,e^{4u}\varepsilon^* + 15\pi e^{2u}\varepsilon^*$) imported from a
lemma stated only for $u \geq 3$, instead of enclosing $\Delta Q$;
the surrogate was numerically true with a factor of about $14$ to spare
but was not the method this appendix describes.
The certificate above encloses $\Delta Q$ directly, as stated.

\noindent\emph{Note on the constant $C=204$.}
The value $C=204$ is a reference value from the original perturbation analysis
(specifically, $|\Delta Q(1)|/(\varepsilon^*(1)\cdot|Q_{\varphi_1}(1)|)$,
not the quotient $|\Delta Q|/(\varepsilon^*\cdot\varphi_1^2)$).
It is \emph{not} used in the certified $B(u)$-bound on $[1,3]$,
and it is \emph{not} the certified far-tail domination constant in
Theorem~\ref{thm:perturbation_tail}, which uses the analytic bound
$2600\,e^{4u}\varepsilon^*(u)$ proved in Appendix~\ref{sec:appE7}.
The primary certified result for the $[1,3]$ region is $B(u) \leq 8.54\times 10^{-24}$
and the sign certificate above, which are independent of the value 204.
It is retained only as a diagnostic value for comparison with earlier
perturbation presentations.
\end{proof}

\begin{lemma}[$\varepsilon^*$ is strictly decreasing]\label{lem:epsstar}
$\varepsilon^*(u) = 2\sum_{n=2}^{\infty}n^4 e^{-\pi(n^2-1)e^{2u}}$ is strictly
decreasing on $[0,\infty)$.
\end{lemma}

\begin{proof}
For each $n\geq 2$:
$\frac{d}{du}[n^4 e^{-\pi(n^2-1)e^{2u}}] = -2\pi(n^2-1)e^{2u}\cdot n^4 e^{-\pi(n^2-1)e^{2u}} < 0$.
The sum of strictly decreasing functions is strictly decreasing.
\end{proof}

\subsection{Far-tail domination for $u\geq 3$}
\label{sec:appE7}

The interval-certified $B(u)$-bound of Lemma~\ref{lem:Bsign} applies only on $[1,3]$.
For the far tail $u\geq 3$ we use a separate explicit domination argument.

\begin{lemma}[Weighted tail comparisons for $u\ge3$]\label{lem:weighted_comparisons}
For $u\ge3$,
\[
\sum_{n\ge2}n^6 e^{-\pi(n^2-1)e^{2u}}
\le
5\sum_{n\ge2}n^4 e^{-\pi(n^2-1)e^{2u}},
\]
and
\[
\sum_{n\ge2}n^8 e^{-\pi(n^2-1)e^{2u}}
\le
32\sum_{n\ge2}n^4 e^{-\pi(n^2-1)e^{2u}}.
\]
\end{lemma}

\begin{proof}
Write $S_m(u) = \sum_{n\ge2}n^m e^{-\pi(n^2-1)e^{2u}}$ and
$a_n = e^{-\pi(n^2-1)e^{2u}}$.
Then $S_6 - 5S_4 = \sum_{n\ge2}(n^6 - 5n^4)a_n$; the $n=2$ term is
$(64 - 80)a_2 = -16a_2$ and every $n\ge3$ term is nonnegative, so
$S_6 \le 5S_4$ if and only if
\[
T_6(u) := \sum_{n\ge3}(n^6 - 5n^4)\,e^{-\pi(n^2-4)e^{2u}} \le 16.
\]
Similarly $S_8 - 32S_4$ has $n=2$ term $(256 - 512)a_2 = -256a_2$, and
$S_8 \le 32S_4$ if and only if
$T_8(u) := \sum_{n\ge3}(n^8 - 32n^4)e^{-\pi(n^2-4)e^{2u}} \le 256$.
Every summand of $T_6$ and $T_8$ is a nonnegative constant times
$e^{-\pi(n^2-4)e^{2u}}$, which is strictly decreasing in~$u$; hence
$T_6$ and $T_8$ are decreasing in~$u$ and the worst case is $u = 3$.
At $u = 3$ the $n=3$ term already dominates each sum, and the Arb/FLINT
certificate \texttt{proof/verify\_weighted\_comparisons.py}
(finite sum $n = 3,\ldots,40$ plus a certified geometric tail bound)
gives
\[
T_6(3) \leq 2.33\times 10^{-2750} < 16,
\qquad
T_8(3) \leq 2.86\times 10^{-2749} < 256,
\]
with margins of over 2{,}700 orders of magnitude.
Thus both comparisons hold for all $u\ge3$.

(Earlier versions of this lemma stated the constants $2$ and $16$;
those are false --- $S_6 > 4S_4$ and $S_8 > 16S_4$ for all $u \ge 3$ by
the same domination ($S_6 - 4S_4$ and $S_8 - 16S_4$ have vanishing
$n=2$ terms and positive $n\ge3$ terms) --- and cited a certificate
script that never evaluated the sums.
The constants $5$ and $32$ proved here are what the downstream
estimates use; see Lemma~\ref{lem:fartail_deriv}.)
\end{proof}

\begin{lemma}[Far-tail derivative-ratio bounds for $u\geq3$]\label{lem:fartail_deriv}
For $u\geq 3$, $n\geq 2$, and $k\in\{0,1,2\}$:
\[
|\varphi_n^{(k)}(u)| \leq a_k n^8 e^{4u} e^{-\pi(n^2-1)e^{2u}} |\varphi_1(u)|,
\]
where $a_0=2$, $a_1=16\pi$, and $a_2=128\pi^2$.
In particular,
\begin{align*}
|R(u)| &\leq \varepsilon^*(u)\,|\varphi_1(u)|,\\
|R'(u)| &\leq 16\pi e^{2u}\varepsilon^*(u)\,|\varphi_1(u)|,\\
|R''(u)| &\leq 128\pi^2 e^{4u}\varepsilon^*(u)\,|\varphi_1(u)|.
\end{align*}
\end{lemma}

\begin{proof}
For $u \geq 3$ the ratio $|g_n/g_1|$ is much smaller than the global
bound of Proposition~\ref{prop:tail_ratio}: since
\[
\frac{g_n}{g_1}
= n^4\,\frac{2\pi^2 e^{9u/2} - 3\pi e^{5u/2}/n^2}
            {2\pi^2 e^{9u/2} - 3\pi e^{5u/2}}
\leq n^4 \bigl(1 - \tfrac{3}{2\pi} e^{-2u}\bigr)^{-1}
\]
and $e^{-2u} \leq e^{-6}$ for $u \geq 3$,
\[
|\varphi_n(u)|/\varphi_1(u) \leq 1.002\, n^4 e^{-\pi(n^2-1)e^{2u}}
\qquad (n \geq 2,\ u \geq 3).
\]
Also $|g_n'/g_n| \leq 5$ and $|g_n''/g_n| \leq 21$ for $n \geq 2$,
$u \geq 3$ (the ratios equal $9\pi^2/(2\pi^2) = 4.5$ and
$(81\pi^2/2)/(2\pi^2) = 20.25$ up to terms of order $e^{-2u}$), so the
relative-derivative formulas of Lemma~\ref{lem:derivbounds} give
\[
|\varphi_n'/\varphi_n| \leq 5 + 2\pi n^2 e^{2u},
\qquad
|\varphi_n''/\varphi_n| \leq 4.01\,\pi^2 n^4 e^{4u},
\]
the second because $21 + 24\pi n^2 e^{2u} \leq 0.01\,\pi^2 n^4 e^{4u}$
for $n \geq 2$, $u \geq 3$.
The per-term display follows with room to spare.
Summing over $n \geq 2$ with
$S_m = \sum_{n\ge2} n^m e^{-\pi(n^2-1)e^{2u}}$ (so $S_4 = \varepsilon^*/2$):

For $k=0$: $|R| \leq 1.002\, S_4\, \varphi_1
= 0.501\,\varepsilon^*\varphi_1 \leq \varepsilon^*\varphi_1$.

For $k=1$: $|R'| \leq 1.002\,(5S_4 + 2\pi e^{2u} S_6)\,\varphi_1
\leq 1.002\,(5 + 10\pi e^{2u})\,S_4\,\varphi_1$
(by Lemma~\ref{lem:weighted_comparisons}, $S_6 \leq 5S_4$)
$= (2.51 + 5.01\pi e^{2u})\,\varepsilon^*\varphi_1
\leq 16\pi e^{2u}\varepsilon^*\varphi_1$ for $u \geq 3$.

For $k=2$: $|R''| \leq 4.02\,\pi^2 e^{4u} S_8\,\varphi_1
\leq 4.02\,\pi^2 e^{4u}\cdot 32 S_4\,\varphi_1$
(by Lemma~\ref{lem:weighted_comparisons}, $S_8 \leq 32S_4$)
$= 64.3\,\pi^2 e^{4u}\varepsilon^*\varphi_1
\leq 128\pi^2 e^{4u}\varepsilon^*\varphi_1$.
\end{proof}

\begin{lemma}[$\Delta Q$ far-tail bound]\label{lem:fartail_deltaQ}
For $u\geq 3$:
\[
\frac{|\Delta Q(u)|}{\varphi_1(u)^2} \leq 2500 e^{4u}\varepsilon^*(u).
\]
\end{lemma}

\begin{proof}
From Lemma~\ref{lem:deltaQ}:
$|\Delta Q|/\varphi_1^2
\leq |R''|/\varphi_1 + (|\varphi_1''|/\varphi_1)(|R|/\varphi_1)
    + 2(|\varphi_1'|/\varphi_1)(|R'|/\varphi_1)
    + (|R''||R| + (R')^2)/\varphi_1^2$.

\emph{Bound on $|\varphi_1'/\varphi_1|$.}
Since
\[
\log\varphi_1=\log\pi+\tfrac{5}{2}u+\log h(u)-\pi e^{2u},
\]
we have
\[
\frac{\varphi_1'}{\varphi_1}
=
\tfrac{5}{2}+\frac{h'}{h}-2\pi e^{2u}.
\]
For $u\ge3$, $h=2\pi e^{2u}-3$, so
\[
0<\frac{h'}{h}
=
\frac{4\pi e^{2u}}{2\pi e^{2u}-3}
<3.
\]
Therefore
\[
\left|\frac{\varphi_1'}{\varphi_1}\right|
\le
\tfrac{5}{2}+3+2\pi e^{2u}
<
3\pi e^{2u}
\]
for $u\ge3$.

\emph{Bound on $|\varphi_1''/\varphi_1|$.}
Using
\[
\frac{\varphi_1''}{\varphi_1}
=
(\log\varphi_1)''
+
\left(\frac{\varphi_1'}{\varphi_1}\right)^{\!2}
=
W_1+
\left(\frac{\varphi_1'}{\varphi_1}\right)^{\!2},
\]
and $|W_1|\le 5\pi e^{2u}$ for $u\ge3$, we get
\[
\left|\frac{\varphi_1''}{\varphi_1}\right|
\le
5\pi e^{2u}+9\pi^2e^{4u}
\le
10\pi^2e^{4u}.
\]

\emph{Linear terms.}
Using Lemma~\ref{lem:fartail_deriv}:
\begin{align*}
|R''|/\varphi_1 &\leq 128\pi^2 e^{4u}\varepsilon^*,\\
(|\varphi_1''|/\varphi_1)(|R|/\varphi_1) &\leq 10\pi^2 e^{4u}\varepsilon^*,\\
2(|\varphi_1'|/\varphi_1)(|R'|/\varphi_1) &\leq 2\cdot3\pi e^{2u}\cdot 16\pi e^{2u}\varepsilon^* = 96\pi^2 e^{4u}\varepsilon^*.
\end{align*}

\emph{Quadratic terms.}
From Lemma~\ref{lem:fartail_deriv},
\[
\frac{|R''||R|}{\varphi_1^2}
\le
128\pi^2 e^{4u}\varepsilon^{*2},
\]
and
\[
\frac{(R')^2}{\varphi_1^2}
\le
(16\pi e^{2u}\varepsilon^*)^2
=
256\pi^2 e^{4u}\varepsilon^{*2}.
\]
Thus
\[
\frac{|R''||R|+(R')^2}{\varphi_1^2}
\le
384\pi^2 e^{4u}\varepsilon^{*2}.
\]
Since $384\pi^2\cdot 10^{-1636}<190$, we have for $u\ge3$
\[
384\pi^2 e^{4u}\varepsilon^{*2}
\le
384\pi^2\cdot 10^{-1636}\cdot e^{4u}\varepsilon^*
<
190\,e^{4u}\varepsilon^*.
\]
Therefore the quadratic terms are absorbed into the gap $2500-2310=190$.

\emph{Conclusion.}
Hence
\[
\frac{|\Delta Q(u)|}{\varphi_1(u)^2}
\le
2500\,e^{4u}\varepsilon^*(u)
\]
for all $u\ge3$.
\end{proof}

\begin{lemma}[$W_{\mathrm{tail}}$ far-tail bound]\label{lem:fartail_Wtail}
For $u\geq 3$:
\[
|W_{\mathrm{tail}}(u)| \leq 2600\,e^{4u}\varepsilon^*(u).
\]
\end{lemma}

\begin{proof}
From Lemma~\ref{lem:Wtail_bound}: $|W_{\mathrm{tail}}| \leq |\Delta Q|/\varphi_1^2 + 3\varepsilon^*|W_1|$.
Lemma~\ref{lem:fartail_deltaQ} gives $|\Delta Q|/\varphi_1^2 \leq 2500e^{4u}\varepsilon^*$.
For the second term: $|W_1(u)| \leq 5\pi e^{2u}$ (Theorem~\ref{thm:algebraic}),
so $3\varepsilon^*|W_1| \leq 15\pi e^{2u}\varepsilon^*$.
For $u\geq3$: $15\pi e^{2u} \leq 15\pi e^{-2u}\cdot e^{4u} < e^{4u}$
(since $15\pi e^{-2u} \leq 15\pi e^{-6} < 1$ for $u\geq3$).
Hence $|W_{\mathrm{tail}}| \leq 2500e^{4u}\varepsilon^* + e^{4u}\varepsilon^* = 2501e^{4u}\varepsilon^*
< 2600\,e^{4u}\varepsilon^*$.
\end{proof}

\noindent\emph{Remark.}
Since $\varepsilon^*(3) < 10^{-1636}$ and $2600e^{12} < 10^7$,
we have $2600e^{12}\varepsilon^*(3) < 10^{-1629} \ll 1$.
Hence the bound in Lemma~\ref{lem:fartail_Wtail} is overwhelmingly small at $u=3$.

\begin{lemma}[Monotonicity of $W_1$ for $u\ge3$]\label{lem:W1_mono}
Let
\[
h(u)=2\pi e^{2u}-3,\qquad
W_1(u)=-\frac{24\pi e^{2u}}{h(u)^2}-4\pi e^{2u}.
\]
Then $W_1'(u)<0$ for all $u\ge3$.
\end{lemma}

\begin{proof}
Differentiate:
\[
W_1'(u)
=
-8\pi e^{2u}
-
48\pi e^{2u}h^{-2}
+
48\pi e^{2u}h^{-3}h'.
\]
Since $h'=4\pi e^{2u}=2(h+3)$, this becomes
\[
W_1'(u)
=
-8\pi e^{2u}
-
48\pi e^{2u}h^{-2}
+
96\pi e^{2u}(h+3)h^{-3}.
\]
Factoring $-8\pi e^{2u}h^{-3}$, we obtain
\[
W_1'(u)
=
-\frac{8\pi e^{2u}}{h^3}
\left(h^3+6h-12(h+3)\right)
=
-\frac{8\pi e^{2u}}{h^3}
\left(h^3-6h-36\right).
\]
For $u\ge3$, $h(u)\geq h(3)=2\pi e^{6}-3>2500$, and
\[
h^3-6h-36>0.
\]
Therefore $W_1'(u)<0$ for all $u\ge3$.
\end{proof}

\begin{lemma}[Monotonicity of far-tail expression]\label{lem:fartail_mono}
The function $u\mapsto e^{4u}\varepsilon^*(u)$ is strictly decreasing on $[3,\infty)$.
Consequently $W_1(u)+2600\,e^{4u}\varepsilon^*(u)$ is strictly decreasing on $[3,\infty)$.
\end{lemma}

\begin{proof}
Since
\[
\varepsilon^*(u)
=
2\sum_{n\ge2}n^4 e^{-\pi(n^2-1)e^{2u}},
\]
differentiation gives
\[
\varepsilon^{*\prime}(u)
=
-4\pi e^{2u}
\sum_{n\ge2} n^4(n^2-1)e^{-\pi(n^2-1)e^{2u}}.
\]
As $n^2-1\ge3$ for $n\ge2$,
\[
\varepsilon^{*\prime}(u)
\le
-12\pi e^{2u}
\sum_{n\ge2}n^4 e^{-\pi(n^2-1)e^{2u}}
=
-6\pi e^{2u}\varepsilon^*(u).
\]
Therefore
\[
\frac{d}{du}\left(e^{4u}\varepsilon^*(u)\right)
=
e^{4u}(4\varepsilon^*(u)+\varepsilon^{*\prime}(u))
\le
e^{4u}(4-6\pi e^{2u})\varepsilon^*(u)<0
\]
for $u\ge3$.
This proves $e^{4u}\varepsilon^*$ is strictly decreasing.
Strict monotonicity of $W_1$ follows from Lemma~\ref{lem:W1_mono}.
\end{proof}

\subsection{Complete proof of Lemma~\ref{lem:wtail}}

\begin{proof}[Proof of Lemma~\ref{lem:wtail}: uniform tail bound]
The proof uses a single route.

\emph{Step 1.} From Lemma~\ref{lem:Wtail_bound}:
$|W_{\mathrm{tail}}(u)| \leq B(u)$
where $B(u) = |\Delta Q(u)|/\varphi_1(u)^2 + 3\varepsilon^*(u)|W_1(u)|$.

\emph{Step 2.} Lemma~\ref{lem:Bsign} certifies by ball arithmetic that
$B(u) \leq 8.54\times 10^{-24}$ for all $u\in[1,3]$
(maximum at the leftmost checkpoint).
In particular $|W_{\mathrm{tail}}(u)| \leq B(u) \leq 8.54\times 10^{-24}$ for all $u\in[1,3]$.

\emph{Step 3.} Lemma~\ref{lem:Bsign} further certifies
$W_1(I)+B(I) < 0$ (an Arb ball comparison)
on each of the 1{,}813 adaptive balls covering $[0.99,3.01]\supset[1,3]$.
Since $W = W_1 + W_{\mathrm{tail}}$ and $|W_{\mathrm{tail}}| \leq B$:
\[
W(u) = W_1(u) + W_{\mathrm{tail}}(u) \leq W_1(u) + B(u) < 0,
\]
with minimum margin $91.29$ at $u\approx1.0$.
\end{proof}

\end{document}
